\documentclass[11pt, ukrainian]{article}
\setlength\textwidth{180mm} \setlength\textheight{277mm}
\setlength\voffset{-38mm} \setlength\hoffset{-28mm}
%\setlength\parindent{0mm}
\pagestyle{empty}
\usepackage[T2A]{fontenc}
\usepackage[utf8]{inputenc}
\usepackage[ukrainian]{babel}
\begin{document}
{\Large \center  Варіанти завдань до лабораторної роботи №4\par 2020/21 н.р.\par \bigskip\bigskip}

Порядок полів у записах та інформацію додаткового файлу фіксовано у варіанті.

Для порівняння числових значень використовується традиційний порядок.

Для порівняння рядкових значень використовується лексикографічний порядок. 

У випадку сортування за кількома полями результуюче сортування будується 
за порядками окремих полів за допомогою конструкції лексикографічного порядку. 
Послідовність полів, за якими має відбуватися сортування, задається в умові. 
Якщо не вказано інше, то поле сортується за зростанням.

Якщо умова вимагає знайти $k$ най… і $k$-е місце займають кілька, то слід виводити всіх їх.

\newpage
\par \bigskip

{\center Варіант  \No \textbf { 1 }\par}
\bigskip
%type = session
%variant = 73
Обробляються результати складання студентами екзаменів зимової сесії.

\underline{Записи основного файлу містять поля}: код групи, оцінка за державною шкалою, предмет, сумарна оцінка в балах за 100-бальною системою, набрані впродовж семестру бали, набрані на екзамені бали, ім’я, прізвище, по-батькові, номер залікової книжки.

\underline{У допоміжному файлі наявні ключі}: кількість предметів, сума балів.

\underline{Знайти} найбільш стабільні предмети. Вивести по кожному з них інформацію:\par
{\noindent}--- на першому рядку:\par
предмет, найменший бал, найбільший бал, стабільність, кількість студентів, що успішно склали;

{\noindent}--- на наступних рядках, починаючи з табуляції, вивести для предмета всіх студентів, що його складали (по одному на рядок):\par
група, середня оцінка по групі в балах (округлена до 2 знаків), кількість студентів групи

{\noindent}у такому сортуванні: група.

%Предмети сортувати: середній бал (округлений)(за спаданням), предмет.

\bigskip\textit{Обмеження}

Студент не може складати предмет більше одного разу.

Бали та оцінка є значеннями цілих числових типів. Решта полів є непорожніми рядками.
Рядки розпочинатись та закінчуватись білими символами не можуть.

Оцінка за державною шкалою може приймати значення: 2-5~--- як зазвичай, 0~--- не з’явився,$-1$~--- не допущений. 
Оцінки 3-5 вважаються задовільними, решта~--- незадовільними. 
Бали є цілими та невід’ємними.
Набрані впродовж семестру бали не можуть перевищувати 60. 
Набрані на екзамені бали не можуть перевищувати 40. 
Набрані на екзамені бали повинні бути не менше 24 або дорівнювати 0(в останньому випадку оцінка не може бути задовільною). 
Сумарна оцінка в балах є сумою балів, набраних у семестрі та на екзамені.
Сумарна оцінка в балах та за державною шкалою мають бути узгоджені між собою.

Прізвище, ім’я та по-батькові (якщо наявне у варіанті) можуть мати до 29, 21, 22 відповідно символів включно кожне. 
Крім літер алфавіту можуть містити апостроф, дефіс та пробіл.

Предмет може мати від 4 до 63 символів включно. 
Крім літер алфавіту може містити подвійні лапки, пробіл та дефіс.

Код групи є непорожнім та складається не більш ніж з 3 літер. 
Крім літер алфавіту може містити десяткові цифри та дефіс.


Номер залікової книжки однозначно ідентифікує студента та складається з 7 десяткових цифр. 

\bigskip\textit{Для деяких варіантів може знадобитись}

Середній бал/оцінка за предметом визначається як середнє арифметичнесумарних балів/державних оцінок за цим предметом. 
Середній бал/оцінка студента визначається як середнє арифметичне сумарних балів/державних оцінокцього студента. 
За обчислення середньої оцінки недопущені та ті, що не з’явилися, рахуються як 2.
У решті випадків недопуски підсумовуються як $-1$, а «не з’явився» як 0.
Якщо студент не гірше ніж задовільно склав усі предмети, 
то його рейтинг за балами/оцінкою визначається як середній бал/оцінка; інакше дорівнює 0.

Стабільність студента визначається як модуль різниці між його кращим та найгіршим сумарним балом.
Стабільність предмета визначається як модуль різниці між найкращим та найгіршим сумарним балом за цим предметом.




\par
\newpage
{\center Варіант  \No \textbf { 2 }\par}
\bigskip
%type = session
%variant = 72
Обробляються результати складання студентами екзаменів зимової сесії.

\underline{Записи основного файлу містять поля}: код групи, предмет, прізвище, сумарна оцінка в балах за 100-бальною системою, набрані на екзамені бали, ім’я, оцінка за державною шкалою, по-батькові, номер залікової книжки.

\underline{У допоміжному файлі наявні ключі}: загальна кількість записів, кількість оцінок 100, SMILE.

\underline{Знайти} предмети з найкращою середньою оцінкою. Вивести по кожному з них інформацію:\par
{\noindent}--- на першому рядку:\par
предмет, середня оцінка (округлений до 1 знаку), кількість студентів, що не склали;

{\noindent}--- на наступних рядках, починаючи з табуляції, вивести для предмета студентів, що отримали з нього менше 95 балів (по одному на рядок):\par
прізвище, ім’я, по-батькові, номер залікової книжки, оцінка в балах, оцінка за державною шкалою

{\noindent}у такому сортуванні: номер залікової книжки.

%Предмети сортувати: кількість студентів, що не склали, середня оцінка (округлена), предмет.

\bigskip\textit{Обмеження}

Студент не може складати предмет більше одного разу.

Бали та оцінка є значеннями цілих числових типів. Решта полів є непорожніми рядками.
Рядки розпочинатись та закінчуватись білими символами не можуть.

Оцінка за державною шкалою може приймати значення: 2-5~--- як зазвичай, 0~--- не з’явився,$-1$~--- не допущений. 
Оцінки 3-5 вважаються задовільними, решта~--- незадовільними. 
Бали є цілими та невід’ємними.
Набрані впродовж семестру бали не можуть перевищувати 60. 
Набрані на екзамені бали не можуть перевищувати 40. 
Набрані на екзамені бали повинні бути не менше 24 або дорівнювати 0(в останньому випадку оцінка не може бути задовільною). 
Сумарна оцінка в балах є сумою балів, набраних у семестрі та на екзамені.
Сумарна оцінка в балах та за державною шкалою мають бути узгоджені між собою.

Прізвище, ім’я та по-батькові (якщо наявне у варіанті) можуть мати до 25, 29, 30 відповідно символів включно кожне. 
Крім літер алфавіту можуть містити апостроф, дефіс та пробіл.

Предмет може мати від 3 до 59 символів включно. 
Крім літер алфавіту може містити подвійні лапки, пробіл та дефіс.

Код групи є непорожнім та складається не більш ніж з 6 літер. 
Крім літер алфавіту може містити десяткові цифри та дефіс.


Номер залікової книжки однозначно ідентифікує студента та складається з 6 десяткових цифр. 

\bigskip\textit{Для деяких варіантів може знадобитись}

Середній бал/оцінка за предметом визначається як середнє арифметичнесумарних балів/державних оцінок за цим предметом. 
Середній бал/оцінка студента визначається як середнє арифметичне сумарних балів/державних оцінокцього студента. 
За обчислення середньої оцінки недопущені та ті, що не з’явилися, рахуються як 2.
У решті випадків недопуски підсумовуються як $-1$, а «не з’явився» як 0.
Якщо студент не гірше ніж задовільно склав усі предмети, 
то його рейтинг за балами/оцінкою визначається як середній бал/оцінка; інакше дорівнює 0.

Стабільність студента визначається як модуль різниці між його кращим та найгіршим сумарним балом.
Стабільність предмета визначається як модуль різниці між найкращим та найгіршимсумарним балом за цим предметом.




\par
\newpage
{\center Варіант  \No \textbf { 3 }\par}
\bigskip
%type = session
%variant = 67
Обробляються результати складання студентами екзаменів зимової сесії.

\underline{Записи основного файлу містять поля}: прізвище, предмет, код групи, набрані на екзамені бали, набрані впродовж семестру бали, оцінка за державною шкалою, ім’я, номер залікової книжки, сумарна оцінка в балах за 100-бальною системою, по-батькові.

\underline{У допоміжному файлі наявні ключі}: кількість студентів без «хвостів», загальна кількість записів.

\underline{Знайти} студентів, що в сесію отримали хоча б по одній оцінці «відмінно», «добре», «задовільно». Вивести по кожному з них інформацію:\par
{\noindent}--- на першому рядку:\par
кількість «добре», ім’я, прізвище, рейтинг за балами (округлений до цілого), номер залікової книжки;

{\noindent}--- на наступних рядках, починаючи з табуляції, вивести для студента всі його результати (по одному на рядок):\par
бали на екзамені, державна оцінка

{\noindent}у такому сортуванні: бали на екзамені, сумарна оцінка, предмет.

%Студентів сортувати: кількість «добре», номер залікової книжки.

\bigskip\textit{Обмеження}

Студент не може складати предмет більше одного разу.

Бали та оцінка є значеннями цілих числових типів. Решта полів є непорожніми рядками.
Рядки розпочинатись та закінчуватись білими символами не можуть.

Оцінка за державною шкалою може приймати значення: 2-5~--- як зазвичай, 0~--- не з’явився,$-1$~--- не допущений. 
Оцінки 3-5 вважаються задовільними, решта~--- незадовільними. 
Бали є цілими та невід’ємними.
Набрані впродовж семестру бали не можуть перевищувати 60. 
Набрані на екзамені бали не можуть перевищувати 40. 
Набрані на екзамені бали повинні бути не менше 24 або дорівнювати 0(в останньому випадку оцінка не може бути задовільною). 
Сумарна оцінка в балах є сумою балів, набраних у семестрі та на екзамені.
Сумарна оцінка в балах та за державною шкалою мають бути узгоджені між собою.

Прізвище, ім’я та по-батькові (якщо наявне у варіанті) можуть мати до 30, 22, 23 відповідно символів включно кожне. 
Крім літер алфавіту можуть містити апостроф, дефіс та пробіл.

Предмет може мати від 5 до 51 символів включно. 
Крім літер алфавіту може містити подвійні лапки, пробіл та дефіс.

Код групи є непорожнім та складається не більш ніж з 7 літер. 
Крім літер алфавіту може містити десяткові цифри та дефіс.


Номер залікової книжки однозначно ідентифікує студента та складається з 8 десяткових цифр. 

\bigskip\textit{Для деяких варіантів може знадобитись}

Середній бал/оцінка за предметом визначається як середнє арифметичнесумарних балів/державних оцінок за цим предметом. 
Середній бал/оцінка студента визначається як середнє арифметичне сумарних балів/державних оцінокцього студента. 
За обчислення середньої оцінки недопущені та ті, що не з’явилися, рахуються як 2.
У решті випадків недопуски підсумовуються як $-1$, а «не з’явився» як 0.
Якщо студент не гірше ніж задовільно склав усі предмети, 
то його рейтинг за балами/оцінкою визначається як середній бал/оцінка; інакше дорівнює 0.

Стабільність студента визначається як модуль різниці між його кращим та найгіршим сумарним балом.
Стабільність предмета визначається як модуль різниці між найкращим та найгіршимсумарним балом за цим предметом.




\par
\newpage
{\center Варіант  \No \textbf { 4 }\par}
\bigskip
%type = session
%variant = 54
Обробляються результати складання студентами екзаменів зимової сесії.

\underline{Записи основного файлу містять поля}: код групи, набрані на екзамені бали, набрані впродовж семестру бали, ім’я, прізвище, предмет, номер залікової книжки, по-батькові, сумарна оцінка в балах за 100-бальною системою, оцінка за державною шкалою.

\underline{У допоміжному файлі наявні ключі}: кількість оцінок «відмінно», кількість студентів, що складали сесію.

\underline{Знайти} 3 найстабільніших студенти. Вивести по кожному з них інформацію:\par
{\noindent}--- на першому рядку:\par
прізвище, ім’я, номер залікової книжки, найменший бал, найбільший бал;

{\noindent}--- на наступних рядках, починаючи з табуляції, вивести для студента всі його результати (по одному на рядок):\par
предмет, сумарна оцінка, державна оцінка прописом, бали на екзамені, бали в семестрі

{\noindent}у такому сортуванні: сумарна оцінка (за спаданням), предмет.

%Студентів сортувати: рейтинг за балами (після округлення), прізвище, ім’я, по-батькові, номер залікової книжки.

\bigskip\textit{Обмеження}

Студент не може складати предмет більше одного разу.

Бали та оцінка є значеннями цілих числових типів. Решта полів є непорожніми рядками.
Рядки розпочинатись та закінчуватись білими символами не можуть.

Оцінка за державною шкалою може приймати значення: 2-5~--- як зазвичай, 0~--- не з’явився,$-1$~--- не допущений. 
Оцінки 3-5 вважаються задовільними, решта~--- незадовільними. 
Бали є цілими та невід’ємними.
Набрані впродовж семестру бали не можуть перевищувати 60. 
Набрані на екзамені бали не можуть перевищувати 40. 
Набрані на екзамені бали повинні бути не менше 24 або дорівнювати 0(в останньому випадку оцінка не може бути задовільною). 
Сумарна оцінка в балах є сумою балів, набраних у семестрі та на екзамені.
Сумарна оцінка в балах та за державною шкалою мають бути узгоджені між собою.

Прізвище, ім’я та по-батькові (якщо наявне у варіанті) можуть мати до 27, 24, 27 відповідно символів включно кожне. 
Крім літер алфавіту можуть містити апостроф, дефіс та пробіл.

Предмет може мати від 3 до 66 символів включно. 
Крім літер алфавіту може містити подвійні лапки, пробіл та дефіс.

Код групи є непорожнім та складається не більш ніж з 4 літер. 
Крім літер алфавіту може містити десяткові цифри та дефіс.


Номер залікової книжки однозначно ідентифікує студента та складається з 7 десяткових цифр. 

\bigskip\textit{Для деяких варіантів може знадобитись}

Середній бал/оцінка за предметом визначається як середнє арифметичнесумарних балів/державних оцінок за цим предметом. 
Середній бал/оцінка студента визначається як середнє арифметичне сумарних балів/державних оцінокцього студента. 
За обчислення середньої оцінки недопущені та ті, що не з’явилися, рахуються як 2.
У решті випадків недопуски підсумовуються як $-1$, а «не з’явився» як 0.
Якщо студент не гірше ніж задовільно склав усі предмети, 
то його рейтинг за балами/оцінкою визначається як середній бал/оцінка; інакше дорівнює 0.

Стабільність студента визначається як модуль різниці між його кращим та найгіршим сумарним балом.
Стабільність предмета визначається як модуль різниці між найкращим та найгіршимсумарним балом за цим предметом.




\par
\newpage
{\center Варіант  \No \textbf { 5 }\par}
\bigskip
%type = session
%variant = 59
Обробляються результати складання студентами екзаменів зимової сесії.

\underline{Записи основного файлу містять поля}: сумарна оцінка в балах за 100-бальною системою, набрані впродовж семестру бали, номер залікової книжки, прізвище, код групи, по-батькові, предмет, ім’я, набрані на екзамені бали, оцінка за державною шкалою.

\underline{У допоміжному файлі наявні ключі}: загальна кількість записів, кількість оцінок 77 балів, кількість оцінок 99 балів.

\underline{Знайти} всіх студентів, що хоча б з одного предмету отримали оцінку 5 за державною шкалою. Вивести по кожному з них інформацію:\par
{\noindent}--- на першому рядку:\par
прізвище, ім’я, по-батькові, номер залікової книжки, кількість «відмінно», кількість «добре», кількість «задовільно», кількість екзаменів;

{\noindent}--- на наступних рядках, починаючи з табуляції, вивести для студента всі його результати (по одному на рядок):\par
сумарна оцінка, предмет, бали на екзамені

{\noindent}у такому сортуванні: бали на екзамені, сумарна оцінка, предмет.

%Студентів сортувати: кількість екзаменів, прізвище, ім’я, по-батькові, номер залікової книжки.

\bigskip\textit{Обмеження}

Студент не може складати предмет більше одного разу.

Бали та оцінка є значеннями цілих числових типів. Решта полів є непорожніми рядками.
Рядки розпочинатись та закінчуватись білими символами не можуть.

Оцінка за державною шкалою може приймати значення: 2-5~--- як зазвичай, 0~--- не з’явився,$-1$~--- не допущений. 
Оцінки 3-5 вважаються задовільними, решта~--- незадовільними. 
Бали є цілими та невід’ємними.
Набрані впродовж семестру бали не можуть перевищувати 60. 
Набрані на екзамені бали не можуть перевищувати 40. 
Набрані на екзамені бали повинні бути не менше 24 або дорівнювати 0(в останньому випадку оцінка не може бути задовільною). 
Сумарна оцінка в балах є сумою балів, набраних у семестрі та на екзамені.
Сумарна оцінка в балах та за державною шкалою мають бути узгоджені між собою.

Прізвище, ім’я та по-батькові (якщо наявне у варіанті) можуть мати до 28, 25, 29 відповідно символів включно кожне. 
Крім літер алфавіту можуть містити апостроф, дефіс та пробіл.

Предмет може мати від 4 до 68 символів включно. 
Крім літер алфавіту може містити подвійні лапки, пробіл та дефіс.

Код групи є непорожнім та складається не більш ніж з 5 літер. 
Крім літер алфавіту може містити десяткові цифри та дефіс.


Номер залікової книжки однозначно ідентифікує студента та складається з 6 десяткових цифр. 

\bigskip\textit{Для деяких варіантів може знадобитись}

Середній бал/оцінка за предметом визначається як середнє арифметичнесумарних балів/державних оцінок за цим предметом. 
Середній бал/оцінка студента визначається як середнє арифметичне сумарних балів/державних оцінокцього студента. 
За обчислення середньої оцінки недопущені та ті, що не з’явилися, рахуються як 2.
У решті випадків недопуски підсумовуються як $-1$, а «не з’явився» як 0.
Якщо студент не гірше ніж задовільно склав усі предмети, 
то його рейтинг за балами/оцінкою визначається як середній бал/оцінка; інакше дорівнює 0.

Стабільність студента визначається як модуль різниці між його кращим та найгіршим сумарним балом.
Стабільність предмета визначається як модуль різниці між найкращим та найгіршимсумарним балом за цим предметом.




\par
\newpage
{\center Варіант  \No \textbf { 6 }\par}
\bigskip
%type = session
%variant = 75
Обробляються результати складання студентами екзаменів зимової сесії.

\underline{Записи основного файлу містять поля}: оцінка за державною шкалою, ім’я, код групи, сумарна оцінка в балах за 100-бальною системою, прізвище, номер залікової книжки, предмет, набрані на екзамені бали, по-батькові.

\underline{У допоміжному файлі наявні ключі}: загальна кількість записів, сума балів.

\underline{Знайти} предмети, середня оцінка за які, вища середньої за всю сесію. Вивести по кожному з них інформацію:\par
{\noindent}--- на першому рядку:\par
середня оцінка (округлена до одного знаку), кількість студентів, що складали, предмет;

{\noindent}--- на наступних рядках, починаючи з табуляції, вивести для предмета всіх студентів, що його складали (по одному на рядок):\par
група, середня оцінка по групі (округлена до цілого), кількість студентів

{\noindent}у такому сортуванні: група.

%Предмети сортувати: середня оцінка (округлена до цілого), предмет.

\bigskip\textit{Обмеження}

Студент не може складати предмет більше одного разу.

Бали та оцінка є значеннями цілих числових типів. Решта полів є непорожніми рядками.
Рядки розпочинатись та закінчуватись білими символами не можуть.

Оцінка за державною шкалою може приймати значення: 2-5~--- як зазвичай, 0~--- не з’явився,$-1$~--- не допущений. 
Оцінки 3-5 вважаються задовільними, решта~--- незадовільними. 
Бали є цілими та невід’ємними.
Набрані впродовж семестру бали не можуть перевищувати 60. 
Набрані на екзамені бали не можуть перевищувати 40. 
Набрані на екзамені бали повинні бути не менше 24 або дорівнювати 0(в останньому випадку оцінка не може бути задовільною). 
Сумарна оцінка в балах є сумою балів, набраних у семестрі та на екзамені.
Сумарна оцінка в балах та за державною шкалою мають бути узгоджені між собою.

Прізвище, ім’я та по-батькові (якщо наявне у варіанті) можуть мати до 22, 20, 21 відповідно символів включно кожне. 
Крім літер алфавіту можуть містити апостроф, дефіс та пробіл.

Предмет може мати від 5 до 56 символів включно. 
Крім літер алфавіту може містити подвійні лапки, пробіл та дефіс.

Код групи є непорожнім та складається не більш ніж з 5 літер. 
Крім літер алфавіту може містити десяткові цифри та дефіс.


Номер залікової книжки однозначно ідентифікує студента та складається з 8 десяткових цифр. 

\bigskip\textit{Для деяких варіантів може знадобитись}

Середній бал/оцінка за предметом визначається як середнє арифметичнесумарних балів/державних оцінок за цим предметом. 
Середній бал/оцінка студента визначається як середнє арифметичне сумарних балів/державних оцінокцього студента. 
За обчислення середньої оцінки недопущені та ті, що не з’явилися, рахуються як 2.
У решті випадків недопуски підсумовуються як $-1$, а «не з’явився» як 0.
Якщо студент не гірше ніж задовільно склав усі предмети, 
то його рейтинг за балами/оцінкою визначається як середній бал/оцінка; інакше дорівнює 0.

Стабільність студента визначається як модуль різниці між його кращим та найгіршим сумарним балом.
Стабільність предмета визначається як модуль різниці між найкращим та найгіршимсумарним балом за цим предметом.




\par
\newpage
{\center Варіант  \No \textbf { 7 }\par}
\bigskip
%type = session
%variant = 57
Обробляються результати складання студентами екзаменів зимової сесії.

\underline{Записи основного файлу містять поля}: сумарна оцінка в балах за 100-бальною системою, набрані впродовж семестру бали, оцінка за державною шкалою, набрані на екзамені бали, код групи, предмет, ім’я, номер залікової книжки, прізвище.

\underline{У допоміжному файлі наявні ключі}: кількість недопусків, кількість студентів без «хвостів», довжина найкоротшого прізвища.

\underline{Знайти} всіх студентів, що по кожному предмету отримали оцінку 5 за державною шкалою. Вивести по кожному з них інформацію:\par
{\noindent}--- на першому рядку:\par
прізвище, ім’я, номер залікової книжки, кількість екзаменів, стабільність, рейтинг за балами (округлений до 1 знаку);

{\noindent}--- на наступних рядках, починаючи з табуляції, вивести для студента всі його результати (по одному на рядок):\par
бали в семестрі, бали на екзамені, сумарна оцінка, предмет

{\noindent}у такому сортуванні: бали в семестрі, предмет.

%Студентів сортувати: кількість екзаменів, рейтинг за балами (після округлення), прізвище, номер залікової книжки.

\bigskip\textit{Обмеження}

Студент не може складати предмет більше одного разу.

Бали та оцінка є значеннями цілих числових типів. Решта полів є непорожніми рядками.
Рядки розпочинатись та закінчуватись білими символами не можуть.

Оцінка за державною шкалою може приймати значення: 2-5~--- як зазвичай, 0~--- не з’явився,$-1$~--- не допущений. 
Оцінки 3-5 вважаються задовільними, решта~--- незадовільними. 
Бали є цілими та невід’ємними.
Набрані впродовж семестру бали не можуть перевищувати 60. 
Набрані на екзамені бали не можуть перевищувати 40. 
Набрані на екзамені бали повинні бути не менше 24 або дорівнювати 0(в останньому випадку оцінка не може бути задовільною). 
Сумарна оцінка в балах є сумою балів, набраних у семестрі та на екзамені.
Сумарна оцінка в балах та за державною шкалою мають бути узгоджені між собою.

Прізвище, ім’я та по-батькові (якщо наявне у варіанті) можуть мати до 23, 23, 24 відповідно символів включно кожне. 
Крім літер алфавіту можуть містити апостроф, дефіс та пробіл.

Предмет може мати від 3 до 69 символів включно. 
Крім літер алфавіту може містити подвійні лапки, пробіл та дефіс.

Код групи є непорожнім та складається не більш ніж з 6 літер. 
Крім літер алфавіту може містити десяткові цифри та дефіс.


Номер залікової книжки однозначно ідентифікує студента та складається з 7 десяткових цифр. 

\bigskip\textit{Для деяких варіантів може знадобитись}

Середній бал/оцінка за предметом визначається як середнє арифметичнесумарних балів/державних оцінок за цим предметом. 
Середній бал/оцінка студента визначається як середнє арифметичне сумарних балів/державних оцінокцього студента. 
За обчислення середньої оцінки недопущені та ті, що не з’явилися, рахуються як 2.
У решті випадків недопуски підсумовуються як $-1$, а «не з’явився» як 0.
Якщо студент не гірше ніж задовільно склав усі предмети, 
то його рейтинг за балами/оцінкою визначається як середній бал/оцінка; інакше дорівнює 0.

Стабільність студента визначається як модуль різниці між його кращим та найгіршим сумарним балом.
Стабільність предмета визначається як модуль різниці між найкращим та найгіршимсумарним балом за цим предметом.




\par
\newpage
{\center Варіант  \No \textbf { 8 }\par}
\bigskip
%type = session
%variant = 55
Обробляються результати складання студентами екзаменів зимової сесії.

\underline{Записи основного файлу містять поля}: предмет, по-батькові, набрані на екзамені бали, оцінка за державною шкалою, прізвище, код групи, набрані впродовж семестру бали, ім’я, номер залікової книжки, сумарна оцінка в балах за 100-бальною системою.

\underline{У допоміжному файлі наявні ключі}: загальна кількість записів, кількість «не з’явився».

\underline{Знайти} 4 найменш стабільних студенти. Вивести по кожному з них інформацію:\par
{\noindent}--- на першому рядку:\par
прізвище, ім’я, по-батькові, номер залікової книжки, стабільність, рейтинг за оцінкою (округлений до 1 знаку);

{\noindent}--- на наступних рядках, починаючи з табуляції, вивести для студента всі його результати (по одному на рядок):\par
державна оцінка, предмет, сумарна оцінка, бали на екзамені, бали в семестрі

{\noindent}у такому сортуванні: державна оцінка, предмет.

%Студентів сортувати: рейтинг за балами (після округлення), прізвище, ім’я, номер залікової книжки.

\bigskip\textit{Обмеження}

Студент не може складати предмет більше одного разу.

Бали та оцінка є значеннями цілих числових типів. Решта полів є непорожніми рядками.
Рядки розпочинатись та закінчуватись білими символами не можуть.

Оцінка за державною шкалою може приймати значення: 2-5~--- як зазвичай, 0~--- не з’явився,$-1$~--- не допущений. 
Оцінки 3-5 вважаються задовільними, решта~--- незадовільними. 
Бали є цілими та невід’ємними.
Набрані впродовж семестру бали не можуть перевищувати 60. 
Набрані на екзамені бали не можуть перевищувати 40. 
Набрані на екзамені бали повинні бути не менше 24 або дорівнювати 0(в останньому випадку оцінка не може бути задовільною). 
Сумарна оцінка в балах є сумою балів, набраних у семестрі та на екзамені.
Сумарна оцінка в балах та за державною шкалою мають бути узгоджені між собою.

Прізвище, ім’я та по-батькові (якщо наявне у варіанті) можуть мати до 21, 26, 26 відповідно символів включно кожне. 
Крім літер алфавіту можуть містити апостроф, дефіс та пробіл.

Предмет може мати від 5 до 54 символів включно. 
Крім літер алфавіту може містити подвійні лапки, пробіл та дефіс.

Код групи є непорожнім та складається не більш ніж з 7 літер. 
Крім літер алфавіту може містити десяткові цифри та дефіс.


Номер залікової книжки однозначно ідентифікує студента та складається з 8 десяткових цифр. 

\bigskip\textit{Для деяких варіантів може знадобитись}

Середній бал/оцінка за предметом визначається як середнє арифметичнесумарних балів/державних оцінок за цим предметом. 
Середній бал/оцінка студента визначається як середнє арифметичне сумарних балів/державних оцінокцього студента. 
За обчислення середньої оцінки недопущені та ті, що не з’явилися, рахуються як 2.
У решті випадків недопуски підсумовуються як $-1$, а «не з’явився» як 0.
Якщо студент не гірше ніж задовільно склав усі предмети, 
то його рейтинг за балами/оцінкою визначається як середній бал/оцінка; інакше дорівнює 0.

Стабільність студента визначається як модуль різниці між його кращим та найгіршим сумарним балом.
Стабільність предмета визначається як модуль різниці між найкращим та найгіршимсумарним балом за цим предметом.




\par
\newpage
{\center Варіант  \No \textbf { 9 }\par}
\bigskip
%type = session
%variant = 58
Обробляються результати складання студентами екзаменів зимової сесії.

\underline{Записи основного файлу містять поля}: оцінка за державною шкалою, ім’я, по-батькові, набрані на екзамені бали, прізвище, сумарна оцінка в балах за 100-бальною системою, предмет, код групи, номер залікової книжки.

\underline{У допоміжному файлі наявні ключі}: кількість студентів, у яких є «хвости», загальна кількість записів.

\underline{Знайти} всіх студентів, що отримали хоча б з одного предмету сумарний бал 100. Вивести по кожному з них інформацію:\par
{\noindent}--- на першому рядку:\par
кількість предметів на 100 балів, номер залікової книжки, прізвище, ім’я, середній бал (округлений до 1 знаку);

{\noindent}--- на наступних рядках, починаючи з табуляції, вивести для студента предмети, що отримали менше 100 балів (по одному на рядок):\par
сумарна оцінка, предмет

{\noindent}у такому сортуванні: сумарна оцінка (за спаданням), предмет.

%Студентів сортувати: кількість предметів на 100 балів (за спаданням), середній бал (після округлення), прізвище, ім’я, по-батькові.

\bigskip\textit{Обмеження}

Студент не може складати предмет більше одного разу.

Бали та оцінка є значеннями цілих числових типів. Решта полів є непорожніми рядками.
Рядки розпочинатись та закінчуватись білими символами не можуть.

Оцінка за державною шкалою може приймати значення: 2-5~--- як зазвичай, 0~--- не з’явився,$-1$~--- не допущений. 
Оцінки 3-5 вважаються задовільними, решта~--- незадовільними. 
Бали є цілими та невід’ємними.
Набрані впродовж семестру бали не можуть перевищувати 60. 
Набрані на екзамені бали не можуть перевищувати 40. 
Набрані на екзамені бали повинні бути не менше 24 або дорівнювати 0(в останньому випадку оцінка не може бути задовільною). 
Сумарна оцінка в балах є сумою балів, набраних у семестрі та на екзамені.
Сумарна оцінка в балах та за державною шкалою мають бути узгоджені між собою.

Прізвище, ім’я та по-батькові (якщо наявне у варіанті) можуть мати до 24, 27, 20 відповідно символів включно кожне. 
Крім літер алфавіту можуть містити апостроф, дефіс та пробіл.

Предмет може мати від 4 до 57 символів включно. 
Крім літер алфавіту може містити подвійні лапки, пробіл та дефіс.

Код групи є непорожнім та складається не більш ніж з 4 літер. 
Крім літер алфавіту може містити десяткові цифри та дефіс.


Номер залікової книжки однозначно ідентифікує студента та складається з 6 десяткових цифр. 

\bigskip\textit{Для деяких варіантів може знадобитись}

Середній бал/оцінка за предметом визначається як середнє арифметичнесумарних балів/державних оцінок за цим предметом. 
Середній бал/оцінка студента визначається як середнє арифметичне сумарних балів/державних оцінокцього студента. 
За обчислення середньої оцінки недопущені та ті, що не з’явилися, рахуються як 2.
У решті випадків недопуски підсумовуються як $-1$, а «не з’явився» як 0.
Якщо студент не гірше ніж задовільно склав усі предмети, 
то його рейтинг за балами/оцінкою визначається як середній бал/оцінка; інакше дорівнює 0.

Стабільність студента визначається як модуль різниці між його кращим та найгіршим сумарним балом.
Стабільність предмета визначається як модуль різниці між найкращим та найгіршимсумарним балом за цим предметом.




\par
\newpage
{\center Варіант  \No \textbf { 10 }\par}
\bigskip
%type = session
%variant = 80
Обробляються результати складання студентами екзаменів зимової сесії.

\underline{Записи основного файлу містять поля}: набрані на екзамені бали, по-батькові, номер залікової книжки, прізвище, оцінка за державною шкалою, предмет, код групи, набрані впродовж семестру бали, сумарна оцінка в балах за 100-бальною системою, ім’я.

\underline{У допоміжному файлі наявні ключі}: сума державних оцінок, загальна кількість предметів.

\underline{Знайти} предмети, які складалися найбільшою кількістю студентів. Вивести по кожному з них інформацію:\par
{\noindent}--- на першому рядку:\par
предмет, середній бал за предметом (округлений до цілого), кількість студентів, що складали;

{\noindent}--- на наступних рядках, починаючи з табуляції, вивести для предмета всіх студентів, що його складали (по одному на рядок):\par
група, середній бал по групі (округлений до 1 знаку), кількість студентів

{\noindent}у такому сортуванні: група.

%Предмети сортувати: стабільність, середній бал (округлений), предмет.

\bigskip\textit{Обмеження}

Студент не може складати предмет більше одного разу.

Бали та оцінка є значеннями цілих числових типів. Решта полів є непорожніми рядками.
Рядки розпочинатись та закінчуватись білими символами не можуть.

Оцінка за державною шкалою може приймати значення: 2-5~--- як зазвичай, 0~--- не з’явився,$-1$~--- не допущений. 
Оцінки 3-5 вважаються задовільними, решта~--- незадовільними. 
Бали є цілими та невід’ємними.
Набрані впродовж семестру бали не можуть перевищувати 60. 
Набрані на екзамені бали не можуть перевищувати 40. 
Набрані на екзамені бали повинні бути не менше 24 або дорівнювати 0(в останньому випадку оцінка не може бути задовільною). 
Сумарна оцінка в балах є сумою балів, набраних у семестрі та на екзамені.
Сумарна оцінка в балах та за державною шкалою мають бути узгоджені між собою.

Прізвище, ім’я та по-батькові (якщо наявне у варіанті) можуть мати до 26, 28, 28 відповідно символів включно кожне. 
Крім літер алфавіту можуть містити апостроф, дефіс та пробіл.

Предмет може мати від 3 до 60 символів включно. 
Крім літер алфавіту може містити подвійні лапки, пробіл та дефіс.

Код групи є непорожнім та складається не більш ніж з 3 літер. 
Крім літер алфавіту може містити десяткові цифри та дефіс.


Номер залікової книжки однозначно ідентифікує студента та складається з 6 десяткових цифр. 

\bigskip\textit{Для деяких варіантів може знадобитись}

Середній бал/оцінка за предметом визначається як середнє арифметичнесумарних балів/державних оцінок за цим предметом. 
Середній бал/оцінка студента визначається як середнє арифметичне сумарних балів/державних оцінокцього студента. 
За обчислення середньої оцінки недопущені та ті, що не з’явилися, рахуються як 2.
У решті випадків недопуски підсумовуються як $-1$, а «не з’явився» як 0.
Якщо студент не гірше ніж задовільно склав усі предмети, 
то його рейтинг за балами/оцінкою визначається як середній бал/оцінка; інакше дорівнює 0.

Стабільність студента визначається як модуль різниці між його кращим та найгіршим сумарним балом.
Стабільність предмета визначається як модуль різниці між найкращим та найгіршимсумарним балом за цим предметом.




\par
\newpage
{\center Варіант  \No \textbf { 11 }\par}
\bigskip
%type = session
%variant = 65
Обробляються результати складання студентами екзаменів зимової сесії.

\underline{Записи основного файлу містять поля}: оцінка за державною шкалою, набрані впродовж семестру бали, набрані на екзамені бали, номер залікової книжки, код групи, прізвище, по-батькові, сумарна оцінка в балах за 100-бальною системою, предмет, ім’я.

\underline{У допоміжному файлі наявні ключі}: кількість записів, скільки разів на екзамені набирали 24 бали.

\underline{Знайти} всіх студентів, чий рейтинг за оцінкою вище середнього рейтингу. Вивести по кожному з них інформацію:\par
{\noindent}--- на першому рядку:\par
ім’я, по-батькові, прізвище, кількість екзаменів, сума балів, номер залікової книжки, рейтинг за оцінкою (округлений до цілого);

{\noindent}--- на наступних рядках, починаючи з табуляції, вивести для студента всі його результати (по одному на рядок):\par
державна оцінка, сумарна оцінка, предмет

{\noindent}у такому сортуванні: сумарний бал, предмет.

%Студентів сортувати: ім’я, прізвище, рейтинг за оцінкою( після округлення), номер залікової книжки.

\bigskip\textit{Обмеження}

Студент не може складати предмет більше одного разу.

Бали та оцінка є значеннями цілих числових типів. Решта полів є непорожніми рядками.
Рядки розпочинатись та закінчуватись білими символами не можуть.

Оцінка за державною шкалою може приймати значення: 2-5~--- як зазвичай, 0~--- не з’явився,$-1$~--- не допущений. 
Оцінки 3-5 вважаються задовільними, решта~--- незадовільними. 
Бали є цілими та невід’ємними.
Набрані впродовж семестру бали не можуть перевищувати 60. 
Набрані на екзамені бали не можуть перевищувати 40. 
Набрані на екзамені бали повинні бути не менше 24 або дорівнювати 0(в останньому випадку оцінка не може бути задовільною). 
Сумарна оцінка в балах є сумою балів, набраних у семестрі та на екзамені.
Сумарна оцінка в балах та за державною шкалою мають бути узгоджені між собою.

Прізвище, ім’я та по-батькові (якщо наявне у варіанті) можуть мати до 20, 30, 25 відповідно символів включно кожне. 
Крім літер алфавіту можуть містити апостроф, дефіс та пробіл.

Предмет може мати від 4 до 50 символів включно. 
Крім літер алфавіту може містити подвійні лапки, пробіл та дефіс.

Код групи є непорожнім та складається не більш ніж з 4 літер. 
Крім літер алфавіту може містити десяткові цифри та дефіс.


Номер залікової книжки однозначно ідентифікує студента та складається з 8 десяткових цифр. 

\bigskip\textit{Для деяких варіантів може знадобитись}

Середній бал/оцінка за предметом визначається як середнє арифметичнесумарних балів/державних оцінок за цим предметом. 
Середній бал/оцінка студента визначається як середнє арифметичне сумарних балів/державних оцінокцього студента. 
За обчислення середньої оцінки недопущені та ті, що не з’явилися, рахуються як 2.
У решті випадків недопуски підсумовуються як $-1$, а «не з’явився» як 0.
Якщо студент не гірше ніж задовільно склав усі предмети, 
то його рейтинг за балами/оцінкою визначається як середній бал/оцінка; інакше дорівнює 0.

Стабільність студента визначається як модуль різниці між його кращим та найгіршим сумарним балом.
Стабільність предмета визначається як модуль різниці між найкращим та найгіршимсумарним балом за цим предметом.




\par
\newpage
{\center Варіант  \No \textbf { 12 }\par}
\bigskip
%type =     session
%variant = 52
Обробляються результати складання студентами екзаменів зимової сесії.

\underline{Записи основного файлу містять поля}: ім’я, номер залікової книжки, код групи, сумарна оцінка в балах за 100-бальною системою, оцінка за державною шкалою, прізвище, предмет, набрані на екзамені бали, по-батькові.

\underline{У допоміжному файлі наявні ключі}: загальна кількість студентів, сума оцінок за державною шкалою.

\underline{Знайти} 5 найкращих студентів (за рейтингом за балами). Вивести по кожному з них інформацію:\par
{\noindent}--- на першому рядку:\par
прізвище, ім’я, по-батькові, номер залікової книжки, рейтинг за балами (округлений до цілого), рейтинг за оцінкою (округлений до 1 знаку);

{\noindent}--- на наступних рядках, починаючи з табуляції, вивести для студента всі його результати (по одному на рядок):\par
предмет, сумарна оцінка, державна оцінка, бали на екзамені, бали в семестрі

{\noindent}у такому сортуванні: сумарна оцінка, предмет.

%Студентів сортувати: рейтинг за балами (після округлення), прізвище, ім’я, по-батькові, номер залікової книжки.

\bigskip\textit{Обмеження}

Студент не може складати предмет більше одного разу.

Бали та оцінка є значеннями цілих числових типів. Решта полів є непорожніми рядками.
Рядки розпочинатись та закінчуватись білими символами не можуть.

Оцінка за державною шкалою може приймати значення: 2-5~--- як зазвичай, 0~--- не з’явився,$-1$~--- не допущений. 
Оцінки 3-5 вважаються задовільними, решта~--- незадовільними. 
Бали є цілими та невід’ємними.
Набрані впродовж семестру бали не можуть перевищувати 60. 
Набрані на екзамені бали не можуть перевищувати 40. 
Набрані на екзамені бали повинні бути не менше 24 або дорівнювати 0(в останньому випадку оцінка не може бути задовільною). 
Сумарна оцінка в балах є сумою балів, набраних у семестрі та на екзамені.
Сумарна оцінка в балах та за державною шкалою мають бути узгоджені між собою.

Прізвище, ім’я та по-батькові (якщо наявне у варіанті) можуть мати до 21, 26, 22 відповідно символів включно кожне. 
Крім літер алфавіту можуть містити апостроф, дефіс та пробіл.

Предмет може мати від 5 до 62 символів включно. 
Крім літер алфавіту може містити подвійні лапки, пробіл та дефіс.

Код групи є непорожнім та складається не більш ніж з 7 літер. 
Крім літер алфавіту може містити десяткові цифри та дефіс.


Номер залікової книжки однозначно ідентифікує студента та складається з 7 десяткових цифр. 

\bigskip\textit{Для деяких варіантів може знадобитись}

Середній бал/оцінка за предметом визначається як середнє арифметичнесумарних балів/державних оцінок за цим предметом. 
Середній бал/оцінка студента визначається як середнє арифметичне сумарних балів/державних оцінокцього студента. 
За обчислення середньої оцінки недопущені та ті, що не з’явилися, рахуються як 2.
У решті випадків недопуски підсумовуються як $-1$, а «не з’явився» як 0.
Якщо студент не гірше ніж задовільно склав усі предмети, 
то його рейтинг за балами/оцінкою визначається як середній бал/оцінка; інакше дорівнює 0.

Стабільність студента визначається як модуль різниці між його кращим та найгіршим сумарним балом.
Стабільність предмета визначається як модуль різниці між найкращим та найгіршимсумарним балом за цим предметом.




\par
\newpage
{\center Варіант  \No \textbf { 13 }\par}
\bigskip
%type = session
%variant = 70
Обробляються результати складання студентами екзаменів зимової сесії.

\underline{Записи основного файлу містять поля}: прізвище, код групи, ім’я, набрані на екзамені бали, предмет, сумарна оцінка в балах за 100-бальною системою, номер залікової книжки, оцінка за державною шкалою, набрані впродовж семестру бали.

\underline{У допоміжному файлі наявні ключі}: загальна кількість предметів, загальна кількість записів.

\underline{Знайти} предмети, за якими кількість незадовільних оцінок найменша. Вивести по кожному з них інформацію:\par
{\noindent}--- на першому рядку:\par
предмет, середній бал (округлений до цілого), кількість студентів, що складали;

{\noindent}--- на наступних рядках, починаючи з табуляції, вивести для предмета всіх студентів, що його складали (по одному на рядок):\par
прізвище, ім’я, по-батькові, номер залікової книжки, група, оцінка в балах, оцінка за державною шкалою, бали за екзамен, бали за семестр

{\noindent}у такому сортуванні: оцінка в балах (за спаданням), номер залікової книжки.

%Предмети сортувати: середній бал (округлений), предмет.

\bigskip\textit{Обмеження}

Студент не може складати предмет більше одного разу.

Бали та оцінка є значеннями цілих числових типів. Решта полів є непорожніми рядками.
Рядки розпочинатись та закінчуватись білими символами не можуть.

Оцінка за державною шкалою може приймати значення: 2-5~--- як зазвичай, 0~--- не з’явився,$-1$~--- не допущений. 
Оцінки 3-5 вважаються задовільними, решта~--- незадовільними. 
Бали є цілими та невід’ємними.
Набрані впродовж семестру бали не можуть перевищувати 60. 
Набрані на екзамені бали не можуть перевищувати 40. 
Набрані на екзамені бали повинні бути не менше 24 або дорівнювати 0(в останньому випадку оцінка не може бути задовільною). 
Сумарна оцінка в балах є сумою балів, набраних у семестрі та на екзамені.
Сумарна оцінка в балах та за державною шкалою мають бути узгоджені між собою.

Прізвище, ім’я та по-батькові (якщо наявне у варіанті) можуть мати до 25, 29, 30 відповідно символів включно кожне. 
Крім літер алфавіту можуть містити апостроф, дефіс та пробіл.

Предмет може мати від 3 до 67 символів включно. 
Крім літер алфавіту може містити подвійні лапки, пробіл та дефіс.

Код групи є непорожнім та складається не більш ніж з 6 літер. 
Крім літер алфавіту може містити десяткові цифри та дефіс.


Номер залікової книжки однозначно ідентифікує студента та складається з 7 десяткових цифр. 

\bigskip\textit{Для деяких варіантів може знадобитись}

Середній бал/оцінка за предметом визначається як середнє арифметичнесумарних балів/державних оцінок за цим предметом. 
Середній бал/оцінка студента визначається як середнє арифметичне сумарних балів/державних оцінокцього студента. 
За обчислення середньої оцінки недопущені та ті, що не з’явилися, рахуються як 2.
У решті випадків недопуски підсумовуються як $-1$, а «не з’явився» як 0.
Якщо студент не гірше ніж задовільно склав усі предмети, 
то його рейтинг за балами/оцінкою визначається як середній бал/оцінка; інакше дорівнює 0.

Стабільність студента визначається як модуль різниці між його кращим та найгіршим сумарним балом.
Стабільність предмета визначається як модуль різниці між найкращим та найгіршимсумарним балом за цим предметом.




\par
\newpage
{\center Варіант  \No \textbf { 14 }\par}
\bigskip
%type = session
%variant = 63
Обробляються результати складання студентами екзаменів зимової сесії.

\underline{Записи основного файлу містять поля}: ім’я, номер залікової книжки, сумарна оцінка в балах за 100-бальною системою, оцінка за державною шкалою, по-батькові, набрані на екзамені бали, предмет, прізвище, код групи.

\underline{У допоміжному файлі наявні ключі}: кількість студентів, у яких були недопуски, загальна кількість записів.

\underline{Знайти} всіх студентів, що по кожному предмету отримали не менше 4 за державною шкалою. Вивести по кожному з них інформацію:\par
{\noindent}--- на першому рядку:\par
середня оцінка (округлена до 2 знаків), кількість екзаменів, номер залікової книжки;

{\noindent}--- на наступних рядках, починаючи з табуляції, вивести для студента всі його результати (по одному на рядок):\par
державна оцінка, предмет, сумарний бал

{\noindent}у такому сортуванні: бали на екзамені, державна оцінка, предмет.

%Студентів сортувати: кількість екзаменів (за спаданням), середня оцінка (після округлення), прізвище, ім’я, по-батькові, номер залікової книжки.

\bigskip\textit{Обмеження}

Студент не може складати предмет більше одного разу.

Бали та оцінка є значеннями цілих числових типів. Решта полів є непорожніми рядками.
Рядки розпочинатись та закінчуватись білими символами не можуть.

Оцінка за державною шкалою може приймати значення: 2-5~--- як зазвичай, 0~--- не з’явився,$-1$~--- не допущений. 
Оцінки 3-5 вважаються задовільними, решта~--- незадовільними. 
Бали є цілими та невід’ємними.
Набрані впродовж семестру бали не можуть перевищувати 60. 
Набрані на екзамені бали не можуть перевищувати 40. 
Набрані на екзамені бали повинні бути не менше 24 або дорівнювати 0(в останньому випадку оцінка не може бути задовільною). 
Сумарна оцінка в балах є сумою балів, набраних у семестрі та на екзамені.
Сумарна оцінка в балах та за державною шкалою мають бути узгоджені між собою.

Прізвище, ім’я та по-батькові (якщо наявне у варіанті) можуть мати до 29, 22, 27 відповідно символів включно кожне. 
Крім літер алфавіту можуть містити апостроф, дефіс та пробіл.

Предмет може мати від 4 до 65 символів включно. 
Крім літер алфавіту може містити подвійні лапки, пробіл та дефіс.

Код групи є непорожнім та складається не більш ніж з 5 літер. 
Крім літер алфавіту може містити десяткові цифри та дефіс.


Номер залікової книжки однозначно ідентифікує студента та складається з 6 десяткових цифр. 

\bigskip\textit{Для деяких варіантів може знадобитись}

Середній бал/оцінка за предметом визначається як середнє арифметичнесумарних балів/державних оцінок за цим предметом. 
Середній бал/оцінка студента визначається як середнє арифметичне сумарних балів/державних оцінокцього студента. 
За обчислення середньої оцінки недопущені та ті, що не з’явилися, рахуються як 2.
У решті випадків недопуски підсумовуються як $-1$, а «не з’явився» як 0.
Якщо студент не гірше ніж задовільно склав усі предмети, 
то його рейтинг за балами/оцінкою визначається як середній бал/оцінка; інакше дорівнює 0.

Стабільність студента визначається як модуль різниці між його кращим та найгіршим сумарним балом.
Стабільність предмета визначається як модуль різниці між найкращим та найгіршимсумарним балом за цим предметом.




\par
\newpage
{\center Варіант  \No \textbf { 15 }\par}
\bigskip
%type = session
%variant = 62
Обробляються результати складання студентами екзаменів зимової сесії.

\underline{Записи основного файлу містять поля}: код групи, ім’я, сумарна оцінка в балах за 100-бальною системою, оцінка за державною шкалою, набрані на екзамені бали, номер залікової книжки, прізвище, предмет, набрані впродовж семестру бали.

\underline{У допоміжному файлі наявні ключі}: загальна кількість записів, сума оцінок.

\underline{Знайти} всіх студентів, що по кожному предмету отримали сумарний бал не менше 75. Вивести по кожному з них інформацію:\par
{\noindent}--- на першому рядку:\par
номер залікової книжки, прізвище, ім’я, стабільність, кількість «відмінно»;

{\noindent}--- на наступних рядках, починаючи з табуляції, вивести для студента всі його результати (по одному на рядок):\par
державна оцінка прописом, предмет

{\noindent}у такому сортуванні: сумарна оцінка, предмет.

%Студентів сортувати: рейтинг за балами (після округлення), номер залікової книжки.

\bigskip\textit{Обмеження}

Студент не може складати предмет більше одного разу.

Бали та оцінка є значеннями цілих числових типів. Решта полів є непорожніми рядками.
Рядки розпочинатись та закінчуватись білими символами не можуть.

Оцінка за державною шкалою може приймати значення: 2-5~--- як зазвичай, 0~--- не з’явився,$-1$~--- не допущений. 
Оцінки 3-5 вважаються задовільними, решта~--- незадовільними. 
Бали є цілими та невід’ємними.
Набрані впродовж семестру бали не можуть перевищувати 60. 
Набрані на екзамені бали не можуть перевищувати 40. 
Набрані на екзамені бали повинні бути не менше 24 або дорівнювати 0(в останньому випадку оцінка не може бути задовільною). 
Сумарна оцінка в балах є сумою балів, набраних у семестрі та на екзамені.
Сумарна оцінка в балах та за державною шкалою мають бути узгоджені між собою.

Прізвище, ім’я та по-батькові (якщо наявне у варіанті) можуть мати до 20, 24, 23 відповідно символів включно кожне. 
Крім літер алфавіту можуть містити апостроф, дефіс та пробіл.

Предмет може мати від 5 до 52 символів включно. 
Крім літер алфавіту може містити подвійні лапки, пробіл та дефіс.

Код групи є непорожнім та складається не більш ніж з 3 літер. 
Крім літер алфавіту може містити десяткові цифри та дефіс.


Номер залікової книжки однозначно ідентифікує студента та складається з 8 десяткових цифр. 

\bigskip\textit{Для деяких варіантів може знадобитись}

Середній бал/оцінка за предметом визначається як середнє арифметичнесумарних балів/державних оцінок за цим предметом. 
Середній бал/оцінка студента визначається як середнє арифметичне сумарних балів/державних оцінокцього студента. 
За обчислення середньої оцінки недопущені та ті, що не з’явилися, рахуються як 2.
У решті випадків недопуски підсумовуються як $-1$, а «не з’явився» як 0.
Якщо студент не гірше ніж задовільно склав усі предмети, 
то його рейтинг за балами/оцінкою визначається як середній бал/оцінка; інакше дорівнює 0.

Стабільність студента визначається як модуль різниці між його кращим та найгіршим сумарним балом.
Стабільність предмета визначається як модуль різниці між найкращим та найгіршимсумарним балом за цим предметом.




\par
\newpage
{\center Варіант  \No \textbf { 16 }\par}
\bigskip
%type = session
%variant = 81
Обробляються результати складання студентами екзаменів зимової сесії.

\underline{Записи основного файлу містять поля}: номер залікової книжки, ім’я, предмет, код групи, по-батько\-ві, сумарна оцінка в балах за 100-бальною системою, оцінка за державною шкалою, прізвище, набрані на екзамені бали.

\underline{У допоміжному файлі наявні ключі}: LISTING, загальна кількість записів, кількість оцінок вище 75 балів.

\underline{Знайти} предмети, які були успішно здані найбільшою кількістю студентів. Вивести по кожному з них інформацію:\par
{\noindent}--- на першому рядку:\par
предмет, середній бал за предметом (округлений до цілого), кількість студентів, що складали, кількість студентів, що склали, кількість «відмінно»;

{\noindent}--- на наступних рядках, починаючи з табуляції, вивести для предмета студентів, що його не склали (по одному на рядок):\par
група, ім’я, прізвище, номер залікової книжки, оцінка в балах, оцінка за державною шкалою

{\noindent}у такому сортуванні: група, прізвище, ім’я, номер залікової книжки.

%Предмети сортувати: кількість «відмінно» (за спаданням), кількість студентів, що складали, предмет.

\bigskip\textit{Обмеження}

Студент не може складати предмет більше одного разу.

Бали та оцінка є значеннями цілих числових типів. Решта полів є непорожніми рядками.
Рядки розпочинатись та закінчуватись білими символами не можуть.

Оцінка за державною шкалою може приймати значення: 2-5~--- як зазвичай, 0~--- не з’явився,$-1$~--- не допущений. 
Оцінки 3-5 вважаються задовільними, решта~--- незадовільними. 
Бали є цілими та невід’ємними.
Набрані впродовж семестру бали не можуть перевищувати 60. 
Набрані на екзамені бали не можуть перевищувати 40. 
Набрані на екзамені бали повинні бути не менше 24 або дорівнювати 0(в останньому випадку оцінка не може бути задовільною). 
Сумарна оцінка в балах є сумою балів, набраних у семестрі та на екзамені.
Сумарна оцінка в балах та за державною шкалою мають бути узгоджені між собою.

Прізвище, ім’я та по-батькові (якщо наявне у варіанті) можуть мати до 24, 23, 25 відповідно символів включно кожне. 
Крім літер алфавіту можуть містити апостроф, дефіс та пробіл.

Предмет може мати від 5 до 61 символів включно. 
Крім літер алфавіту може містити подвійні лапки, пробіл та дефіс.

Код групи є непорожнім та складається не більш ніж з 5 літер. 
Крім літер алфавіту може містити десяткові цифри та дефіс.


Номер залікової книжки однозначно ідентифікує студента та складається з 7 десяткових цифр. 

\bigskip\textit{Для деяких варіантів може знадобитись}

Середній бал/оцінка за предметом визначається як середнє арифметичнесумарних балів/державних оцінок за цим предметом. 
Середній бал/оцінка студента визначається як середнє арифметичне сумарних балів/державних оцінокцього студента. 
За обчислення середньої оцінки недопущені та ті, що не з’явилися, рахуються як 2.
У решті випадків недопуски підсумовуються як $-1$, а «не з’явився» як 0.
Якщо студент не гірше ніж задовільно склав усі предмети, 
то його рейтинг за балами/оцінкою визначається як середній бал/оцінка; інакше дорівнює 0.

Стабільність студента визначається як модуль різниці між його кращим та найгіршим сумарним балом.
Стабільність предмета визначається як модуль різниці між найкращим та найгіршимсумарним балом за цим предметом.




\par
\newpage
{\center Варіант  \No \textbf { 17 }\par}
\bigskip
%type = session
%variant = 53
Обробляються результати складання студентами екзаменів зимової сесії.

\underline{Записи основного файлу містять поля}: код групи, набрані на екзамені бали, набрані впродовж семестру бали, сумарна оцінка в балах за 100-бальною системою, предмет, ім’я, оцінка за державною шкалою, прізвище, номер залікової книжки.

\underline{У допоміжному файлі наявні ключі}: найбільша довжина предмету, загальна кількість записів.

\underline{Знайти} 7 найкращих студентів (за рейтингом за національною шкалою). Вивести по кожному з них інформацію:\par
{\noindent}--- на першому рядку:\par
прізвище, ім’я, номер залікової книжки, рейтинг за балами (округлений до цілого), рейтинг за оцінкою (округлений до 2 знаків);

{\noindent}--- на наступних рядках, починаючи з табуляції, вивести для студента всі предмети, де було отримано не «відмінно» (по одному на рядок):\par
сумарна оцінка, бали в семестрі, державна оцінка, предмет

{\noindent}у такому сортуванні: предмет.

%Студентів сортувати: прізвище, ім’я, сума балів за сесію, номер залікової книжки.

\bigskip\textit{Обмеження}

Студент не може складати предмет більше одного разу.

Бали та оцінка є значеннями цілих числових типів. Решта полів є непорожніми рядками.
Рядки розпочинатись та закінчуватись білими символами не можуть.

Оцінка за державною шкалою може приймати значення: 2-5~--- як зазвичай, 0~--- не з’явився,$-1$~--- не допущений. 
Оцінки 3-5 вважаються задовільними, решта~--- незадовільними. 
Бали є цілими та невід’ємними.
Набрані впродовж семестру бали не можуть перевищувати 60. 
Набрані на екзамені бали не можуть перевищувати 40. 
Набрані на екзамені бали повинні бути не менше 24 або дорівнювати 0(в останньому випадку оцінка не може бути задовільною). 
Сумарна оцінка в балах є сумою балів, набраних у семестрі та на екзамені.
Сумарна оцінка в балах та за державною шкалою мають бути узгоджені між собою.

Прізвище, ім’я та по-батькові (якщо наявне у варіанті) можуть мати до 23, 30, 21 відповідно символів включно кожне. 
Крім літер алфавіту можуть містити апостроф, дефіс та пробіл.

Предмет може мати від 3 до 58 символів включно. 
Крім літер алфавіту може містити подвійні лапки, пробіл та дефіс.

Код групи є непорожнім та складається не більш ніж з 7 літер. 
Крім літер алфавіту може містити десяткові цифри та дефіс.


Номер залікової книжки однозначно ідентифікує студента та складається з 6 десяткових цифр. 

\bigskip\textit{Для деяких варіантів може знадобитись}

Середній бал/оцінка за предметом визначається як середнє арифметичнесумарних балів/державних оцінок за цим предметом. 
Середній бал/оцінка студента визначається як середнє арифметичне сумарних балів/державних оцінокцього студента. 
За обчислення середньої оцінки недопущені та ті, що не з’явилися, рахуються як 2.
У решті випадків недопуски підсумовуються як $-1$, а «не з’явився» як 0.
Якщо студент не гірше ніж задовільно склав усі предмети, 
то його рейтинг за балами/оцінкою визначається як середній бал/оцінка; інакше дорівнює 0.

Стабільність студента визначається як модуль різниці між його кращим та найгіршим сумарним балом.
Стабільність предмета визначається як модуль різниці між найкращим та найгіршимсумарним балом за цим предметом.




\par
\newpage
{\center Варіант  \No \textbf { 18 }\par}
\bigskip
%type = session
%variant = 71
Обробляються результати складання студентами екзаменів зимової сесії.

\underline{Записи основного файлу містять поля}: номер залікової книжки, предмет, код групи, сумарна оцінка в балах за 100-бальною системою, набрані на екзамені бали, набрані впродовж семестру бали, прізвище, ім’я, оцінка за державною шкалою, по-батькові.

\underline{У допоміжному файлі наявні ключі}: довжина найдовшого прізвища, загальна кількість «хвостів».

\underline{Знайти} предмети з найкращим середнім балом. Вивести по кожному з них інформацію:\par
{\noindent}--- на першому рядку:\par
предмет, середній бал (округлений до цілого), середня оцінка (округлена до 1 знаку);

{\noindent}--- на наступних рядках, починаючи з табуляції, вивести для предмета всіх студентів, що його складали (по одному на рядок):\par
оцінка в балах, бали за екзамен, прізвище, ім’я, група, номер залікової книжки

{\noindent}у такому сортуванні: прізвище, ім’я, номер залікової книжки.

%Предмети сортувати: середній бал (округлений), предмет.

\bigskip\textit{Обмеження}

Студент не може складати предмет більше одного разу.

Бали та оцінка є значеннями цілих числових типів. Решта полів є непорожніми рядками.
Рядки розпочинатись та закінчуватись білими символами не можуть.

Оцінка за державною шкалою може приймати значення: 2-5~--- як зазвичай, 0~--- не з’явився,$-1$~--- не допущений. 
Оцінки 3-5 вважаються задовільними, решта~--- незадовільними. 
Бали є цілими та невід’ємними.
Набрані впродовж семестру бали не можуть перевищувати 60. 
Набрані на екзамені бали не можуть перевищувати 40. 
Набрані на екзамені бали повинні бути не менше 24 або дорівнювати 0(в останньому випадку оцінка не може бути задовільною). 
Сумарна оцінка в балах є сумою балів, набраних у семестрі та на екзамені.
Сумарна оцінка в балах та за державною шкалою мають бути узгоджені між собою.

Прізвище, ім’я та по-батькові (якщо наявне у варіанті) можуть мати до 22, 25, 20 відповідно символів включно кожне. 
Крім літер алфавіту можуть містити апостроф, дефіс та пробіл.

Предмет може мати від 4 до 55 символів включно. 
Крім літер алфавіту може містити подвійні лапки, пробіл та дефіс.

Код групи є непорожнім та складається не більш ніж з 6 літер. 
Крім літер алфавіту може містити десяткові цифри та дефіс.


Номер залікової книжки однозначно ідентифікує студента та складається з 8 десяткових цифр. 

\bigskip\textit{Для деяких варіантів може знадобитись}

Середній бал/оцінка за предметом визначається як середнє арифметичнесумарних балів/державних оцінок за цим предметом. 
Середній бал/оцінка студента визначається як середнє арифметичне сумарних балів/державних оцінокцього студента. 
За обчислення середньої оцінки недопущені та ті, що не з’явилися, рахуються як 2.
У решті випадків недопуски підсумовуються як $-1$, а «не з’явився» як 0.
Якщо студент не гірше ніж задовільно склав усі предмети, 
то його рейтинг за балами/оцінкою визначається як середній бал/оцінка; інакше дорівнює 0.

Стабільність студента визначається як модуль різниці між його кращим та найгіршим сумарним балом.
Стабільність предмета визначається як модуль різниці між найкращим та найгіршимсумарним балом за цим предметом.




\par
\newpage
{\center Варіант  \No \textbf { 19 }\par}
\bigskip
%type = session
%variant = 68
Обробляються результати складання студентами екзаменів зимової сесії.

\underline{Записи основного файлу містять поля}: набрані на екзамені бали, оцінка за державною шкалою, ім’я, предмет, набрані впродовж семестру бали, сумарна оцінка в балах за 100-бальною системою, код групи, прізвище, номер залікової книжки.

\underline{У допоміжному файлі наявні ключі}: сума державних оцінок, кількість записів.

\underline{Знайти} всіх студентів, що не склали жодного предмета. Вивести по кожному з них інформацію:\par
{\noindent}--- на першому рядку:\par
ім’я, прізвище, номер залікової книжки, кількість «хвостів», кількість «незадовільно»;

{\noindent}--- на наступних рядках, починаючи з табуляції, вивести для студента всі його результати (по одному на рядок):\par
предмет, бали в семестрі, бали на екзамені, сумарна оцінка, державна оцінка

{\noindent}у такому сортуванні: предмет.

%Студентів сортувати: кількість «хвостів» (за спаданням), прізвище, номер залікової книжки.

\bigskip\textit{Обмеження}

Студент не може складати предмет більше одного разу.

Бали та оцінка є значеннями цілих числових типів. Решта полів є непорожніми рядками.
Рядки розпочинатись та закінчуватись білими символами не можуть.

Оцінка за державною шкалою може приймати значення: 2-5~--- як зазвичай, 0~--- не з’явився,$-1$~--- не допущений. 
Оцінки 3-5 вважаються задовільними, решта~--- незадовільними. 
Бали є цілими та невід’ємними.
Набрані впродовж семестру бали не можуть перевищувати 60. 
Набрані на екзамені бали не можуть перевищувати 40. 
Набрані на екзамені бали повинні бути не менше 24 або дорівнювати 0(в останньому випадку оцінка не може бути задовільною). 
Сумарна оцінка в балах є сумою балів, набраних у семестрі та на екзамені.
Сумарна оцінка в балах та за державною шкалою мають бути узгоджені між собою.

Прізвище, ім’я та по-батькові (якщо наявне у варіанті) можуть мати до 27, 28, 26 відповідно символів включно кожне. 
Крім літер алфавіту можуть містити апостроф, дефіс та пробіл.

Предмет може мати від 4 до 53 символів включно. 
Крім літер алфавіту може містити подвійні лапки, пробіл та дефіс.

Код групи є непорожнім та складається не більш ніж з 4 літер. 
Крім літер алфавіту може містити десяткові цифри та дефіс.


Номер залікової книжки однозначно ідентифікує студента та складається з 8 десяткових цифр. 

\bigskip\textit{Для деяких варіантів може знадобитись}

Середній бал/оцінка за предметом визначається як середнє арифметичнесумарних балів/державних оцінок за цим предметом. 
Середній бал/оцінка студента визначається як середнє арифметичне сумарних балів/державних оцінокцього студента. 
За обчислення середньої оцінки недопущені та ті, що не з’явилися, рахуються як 2.
У решті випадків недопуски підсумовуються як $-1$, а «не з’явився» як 0.
Якщо студент не гірше ніж задовільно склав усі предмети, 
то його рейтинг за балами/оцінкою визначається як середній бал/оцінка; інакше дорівнює 0.

Стабільність студента визначається як модуль різниці між його кращим та найгіршим сумарним балом.
Стабільність предмета визначається як модуль різниці між найкращим та найгіршимсумарним балом за цим предметом.




\par
\newpage
{\center Варіант  \No \textbf { 20 }\par}
\bigskip
%type = session
%variant = 78
Обробляються результати складання студентами екзаменів зимової сесії.

\underline{Записи основного файлу містять поля}: сумарна оцінка в балах за 100-бальною системою, номер залікової книжки, набрані впродовж семестру бали, код групи, прізвище, предмет, оцінка за державною шкалою, ім’я, набрані на екзамені бали.

\underline{У допоміжному файлі наявні ключі}: загальна кількість предметів, кількість «не з’явився».

\underline{Знайти} предмети, які були складені менше, ніж половиною студентів. Вивести по кожному з них інформацію:\par
{\noindent}--- на першому рядку:\par
кількість, що складали, кількість, що не склали, середня оцінка (округлена до 1 знаку), предмет;

{\noindent}--- на наступних рядках, починаючи з табуляції, вивести для предмета студентів, що його успішно склали (по одному на рядок):\par
ім’я, прізвище, номер залікової книжки, оцінка в балах, оцінка за державною шкалою, бали за екзамен, бали за семестр

{\noindent}у такому сортуванні: прізвище, ім’я, номер залікової книжки.

%Предмети сортувати: середня оцінка (округлена), предмет.

\bigskip\textit{Обмеження}

Студент не може складати предмет більше одного разу.

Бали та оцінка є значеннями цілих числових типів. Решта полів є непорожніми рядками.
Рядки розпочинатись та закінчуватись білими символами не можуть.

Оцінка за державною шкалою може приймати значення: 2-5~--- як зазвичай, 0~--- не з’явився,$-1$~--- не допущений. 
Оцінки 3-5 вважаються задовільними, решта~--- незадовільними. 
Бали є цілими та невід’ємними.
Набрані впродовж семестру бали не можуть перевищувати 60. 
Набрані на екзамені бали не можуть перевищувати 40. 
Набрані на екзамені бали повинні бути не менше 24 або дорівнювати 0(в останньому випадку оцінка не може бути задовільною). 
Сумарна оцінка в балах є сумою балів, набраних у семестрі та на екзамені.
Сумарна оцінка в балах та за державною шкалою мають бути узгоджені між собою.

Прізвище, ім’я та по-батькові (якщо наявне у варіанті) можуть мати до 26, 27, 29 відповідно символів включно кожне. 
Крім літер алфавіту можуть містити апостроф, дефіс та пробіл.

Предмет може мати від 5 до 70 символів включно. 
Крім літер алфавіту може містити подвійні лапки, пробіл та дефіс.

Код групи є непорожнім та складається не більш ніж з 3 літер. 
Крім літер алфавіту може містити десяткові цифри та дефіс.


Номер залікової книжки однозначно ідентифікує студента та складається з 6 десяткових цифр. 

\bigskip\textit{Для деяких варіантів може знадобитись}

Середній бал/оцінка за предметом визначається як середнє арифметичнесумарних балів/державних оцінок за цим предметом. 
Середній бал/оцінка студента визначається як середнє арифметичне сумарних балів/державних оцінокцього студента. 
За обчислення середньої оцінки недопущені та ті, що не з’явилися, рахуються як 2.
У решті випадків недопуски підсумовуються як $-1$, а «не з’явився» як 0.
Якщо студент не гірше ніж задовільно склав усі предмети, 
то його рейтинг за балами/оцінкою визначається як середній бал/оцінка; інакше дорівнює 0.

Стабільність студента визначається як модуль різниці між його кращим та найгіршим сумарним балом.
Стабільність предмета визначається як модуль різниці між найкращим та найгіршимсумарним балом за цим предметом.




\par
\newpage
{\center Варіант  \No \textbf { 21 }\par}
\bigskip
%type = session
%variant = 60
Обробляються результати складання студентами екзаменів зимової сесії.

\underline{Записи основного файлу містять поля}: номер залікової книжки, набрані впродовж семестру бали, сумарна оцінка в балах за 100-бальною системою, набрані на екзамені бали, прізвище, ім’я, код групи, предмет, оцінка за державною шкалою.

\underline{У допоміжному файлі наявні ключі}: довжина найкоротшого прізвища, загальна кількість записів.

\underline{Знайти} всіх студентів, що по кожному предмету отримали сумарний бал рівно 60. Вивести по кожному з них інформацію:\par
{\noindent}--- на першому рядку:\par
номер залікової книжки, прізвище, ім’я, найменший бал на екзамені, середній бал в семестрі (округлений до цілого);

{\noindent}--- на наступних рядках, починаючи з табуляції, вивести для студента всі його результати (по одному на рядок):\par
предмет, бали на екзамені, бали в семестрі

{\noindent}у такому сортуванні: бали в семестрі, предмет.

%Студентів сортувати: середній бал в семестрі (після округлення), номер залікової книжки.

\bigskip\textit{Обмеження}

Студент не може складати предмет більше одного разу.

Бали та оцінка є значеннями цілих числових типів. Решта полів є непорожніми рядками.
Рядки розпочинатись та закінчуватись білими символами не можуть.

Оцінка за державною шкалою може приймати значення: 2-5~--- як зазвичай, 0~--- не з’явився,$-1$~--- не допущений. 
Оцінки 3-5 вважаються задовільними, решта~--- незадовільними. 
Бали є цілими та невід’ємними.
Набрані впродовж семестру бали не можуть перевищувати 60. 
Набрані на екзамені бали не можуть перевищувати 40. 
Набрані на екзамені бали повинні бути не менше 24 або дорівнювати 0(в останньому випадку оцінка не може бути задовільною). 
Сумарна оцінка в балах є сумою балів, набраних у семестрі та на екзамені.
Сумарна оцінка в балах та за державною шкалою мають бути узгоджені між собою.

Прізвище, ім’я та по-батькові (якщо наявне у варіанті) можуть мати до 28, 20, 24 відповідно символів включно кожне. 
Крім літер алфавіту можуть містити апостроф, дефіс та пробіл.

Предмет може мати від 3 до 64 символів включно. 
Крім літер алфавіту може містити подвійні лапки, пробіл та дефіс.

Код групи є непорожнім та складається не більш ніж з 4 літер. 
Крім літер алфавіту може містити десяткові цифри та дефіс.


Номер залікової книжки однозначно ідентифікує студента та складається з 7 десяткових цифр. 

\bigskip\textit{Для деяких варіантів може знадобитись}

Середній бал/оцінка за предметом визначається як середнє арифметичнесумарних балів/державних оцінок за цим предметом. 
Середній бал/оцінка студента визначається як середнє арифметичне сумарних балів/державних оцінокцього студента. 
За обчислення середньої оцінки недопущені та ті, що не з’явилися, рахуються як 2.
У решті випадків недопуски підсумовуються як $-1$, а «не з’явився» як 0.
Якщо студент не гірше ніж задовільно склав усі предмети, 
то його рейтинг за балами/оцінкою визначається як середній бал/оцінка; інакше дорівнює 0.

Стабільність студента визначається як модуль різниці між його кращим та найгіршим сумарним балом.
Стабільність предмета визначається як модуль різниці між найкращим та найгіршимсумарним балом за цим предметом.




\par
\newpage
{\center Варіант  \No \textbf { 22 }\par}
\bigskip
%type = session
%variant = 69
Обробляються результати складання студентами екзаменів зимової сесії.

\underline{Записи основного файлу містять поля}: номер залікової книжки, оцінка за державною шкалою, предмет, по-батькові, набрані на екзамені бали, прізвище, ім’я, сумарна оцінка в балах за 100-бальною системою, код групи.

\underline{У допоміжному файлі наявні ключі}: кількість «відмінно», сума державних оцінок.

\underline{Знайти} всіх студентів, що щось не склали і при цьому отримували оцінку «відмінно». Вивести по кожному з них інформацію:\par
{\noindent}--- на першому рядку:\par
кількість «відмінно», кількість «хвостів», ім’я, прізвище, рейтинг за оцінкою (округлений до 1 знаку), номер залікової книжки;

{\noindent}--- на наступних рядках, починаючи з табуляції, вивести для студента предмети, з оцінкою нижче «відмінно» (по одному на рядок):\par
предмет, сумарна оцінка, державна оцінка

{\noindent}у такому сортуванні: предмет, сумарна оцінка.

%Студентів сортувати: кількість «відмінно», кількість «хвостів» (за спаданням), прізвище, номер залікової книжки.

\bigskip\textit{Обмеження}

Студент не може складати предмет більше одного разу.

Бали та оцінка є значеннями цілих числових типів. Решта полів є непорожніми рядками.
Рядки розпочинатись та закінчуватись білими символами не можуть.

Оцінка за державною шкалою може приймати значення: 2-5~--- як зазвичай, 0~--- не з’явився,$-1$~--- не допущений. 
Оцінки 3-5 вважаються задовільними, решта~--- незадовільними. 
Бали є цілими та невід’ємними.
Набрані впродовж семестру бали не можуть перевищувати 60. 
Набрані на екзамені бали не можуть перевищувати 40. 
Набрані на екзамені бали повинні бути не менше 24 або дорівнювати 0(в останньому випадку оцінка не може бути задовільною). 
Сумарна оцінка в балах є сумою балів, набраних у семестрі та на екзамені.
Сумарна оцінка в балах та за державною шкалою мають бути узгоджені між собою.

Прізвище, ім’я та по-батькові (якщо наявне у варіанті) можуть мати до 30, 21, 28 відповідно символів включно кожне. 
Крім літер алфавіту можуть містити апостроф, дефіс та пробіл.

Предмет може мати від 3 до 58 символів включно. 
Крім літер алфавіту може містити подвійні лапки, пробіл та дефіс.

Код групи є непорожнім та складається не більш ніж з 5 літер. 
Крім літер алфавіту може містити десяткові цифри та дефіс.


Номер залікової книжки однозначно ідентифікує студента та складається з 7 десяткових цифр. 

\bigskip\textit{Для деяких варіантів може знадобитись}

Середній бал/оцінка за предметом визначається як середнє арифметичнесумарних балів/державних оцінок за цим предметом. 
Середній бал/оцінка студента визначається як середнє арифметичне сумарних балів/державних оцінокцього студента. 
За обчислення середньої оцінки недопущені та ті, що не з’явилися, рахуються як 2.
У решті випадків недопуски підсумовуються як $-1$, а «не з’явився» як 0.
Якщо студент не гірше ніж задовільно склав усі предмети, 
то його рейтинг за балами/оцінкою визначається як середній бал/оцінка; інакше дорівнює 0.

Стабільність студента визначається як модуль різниці між його кращим та найгіршим сумарним балом.
Стабільність предмета визначається як модуль різниці між найкращим та найгіршимсумарним балом за цим предметом.




\par
\newpage
{\center Варіант  \No \textbf { 23 }\par}
\bigskip
%type = session
%variant = 66
Обробляються результати складання студентами екзаменів зимової сесії.

\underline{Записи основного файлу містять поля}: предмет, номер залікової книжки, ім’я, набрані на екзамені бали, код групи, по-батькові, прізвище, оцінка за державною шкалою, сумарна оцінка в балах за 100-бальною системою.

\underline{У допоміжному файлі наявні ключі}: найбільша довжина назви предмета, найменша довжина назви предмета, найбільший номер залікової книжки.

\underline{Знайти} студентів, що не отримали жодної оцінки «добре». Вивести по кожному з них інформацію:\par
{\noindent}--- на першому рядку:\par
кількість «хвостів», кількість екзаменів, номер залікової книжки, прізвище, ім’я, по-батькові;

{\noindent}--- на наступних рядках, починаючи з табуляції, вивести для студента всі його результати (по одному на рядок):\par
бали в семестрі, бали на екзамені, предмет, державна оцінка

{\noindent}у такому сортуванні: бали на екзамені, предмет.

%Студентів сортувати: кількість «хвостів», ім’я, прізвище, номер залікової книжки.

\bigskip\textit{Обмеження}

Студент не може складати предмет більше одного разу.

Бали та оцінка є значеннями цілих числових типів. Решта полів є непорожніми рядками.
Рядки розпочинатись та закінчуватись білими символами не можуть.

Оцінка за державною шкалою може приймати значення: 2-5~--- як зазвичай, 0~--- не з’явився,$-1$~--- не допущений. 
Оцінки 3-5 вважаються задовільними, решта~--- незадовільними. 
Бали є цілими та невід’ємними.
Набрані впродовж семестру бали не можуть перевищувати 60. 
Набрані на екзамені бали не можуть перевищувати 40. 
Набрані на екзамені бали повинні бути не менше 24 або дорівнювати 0(в останньому випадку оцінка не може бути задовільною). 
Сумарна оцінка в балах є сумою балів, набраних у семестрі та на екзамені.
Сумарна оцінка в балах та за державною шкалою мають бути узгоджені між собою.

Прізвище, ім’я та по-батькові (якщо наявне у варіанті) можуть мати до 22, 21, 29 відповідно символів включно кожне. 
Крім літер алфавіту можуть містити апостроф, дефіс та пробіл.

Предмет може мати від 5 до 57 символів включно. 
Крім літер алфавіту може містити подвійні лапки, пробіл та дефіс.

Код групи є непорожнім та складається не більш ніж з 6 літер. 
Крім літер алфавіту може містити десяткові цифри та дефіс.


Номер залікової книжки однозначно ідентифікує студента та складається з 6 десяткових цифр. 

\bigskip\textit{Для деяких варіантів може знадобитись}

Середній бал/оцінка за предметом визначається як середнє арифметичнесумарних балів/державних оцінок за цим предметом. 
Середній бал/оцінка студента визначається як середнє арифметичне сумарних балів/державних оцінокцього студента. 
За обчислення середньої оцінки недопущені та ті, що не з’явилися, рахуються як 2.
У решті випадків недопуски підсумовуються як $-1$, а «не з’явився» як 0.
Якщо студент не гірше ніж задовільно склав усі предмети, 
то його рейтинг за балами/оцінкою визначається як середній бал/оцінка; інакше дорівнює 0.

Стабільність студента визначається як модуль різниці між його кращим та найгіршим сумарним балом.
Стабільність предмета визначається як модуль різниці між найкращим та найгіршимсумарним балом за цим предметом.




\par
\newpage
{\center Варіант  \No \textbf { 24 }\par}
\bigskip
%type = session
%variant = 79
Обробляються результати складання студентами екзаменів зимової сесії.

\underline{Записи основного файлу містять поля}: прізвище, набрані на екзамені бали, предмет, номер залікової книжки, оцінка за державною шкалою, ім’я, код групи, сумарна оцінка в балах за 100-бальною системою, набрані впродовж семестру бали.

\underline{У допоміжному файлі наявні ключі}: FAILED SUCCESSFULLY, загальна кількість записів, найменший номер залікової книжки.

\underline{Знайти} предмети, за якими оцінок «відмінно» було більше, ніж оцінок «задовільно». Вивести по кожному з них інформацію:\par
{\noindent}--- на першому рядку:\par
кількість «відмінно» -(мінус) кількість «задовільно», предмет;

{\noindent}--- на наступних рядках, починаючи з табуляції, вивести для предмета тільки тих, хто склав на «відмінно» або «задовільно» (по одному на рядок):\par
ім’я, прізвище, номер залікової книжки, група, оцінка в балах

{\noindent}у такому сортуванні: прізвище, ім’я, номер залікової книжки.

%Предмети сортувати: за першим полем, що виводиться, предмет.

\bigskip\textit{Обмеження}

Студент не може складати предмет більше одного разу.

Бали та оцінка є значеннями цілих числових типів. Решта полів є непорожніми рядками.
Рядки розпочинатись та закінчуватись білими символами не можуть.

Оцінка за державною шкалою може приймати значення: 2-5~--- як зазвичай, 0~--- не з’явився,$-1$~--- не допущений. 
Оцінки 3-5 вважаються задовільними, решта~--- незадовільними. 
Бали є цілими та невід’ємними.
Набрані впродовж семестру бали не можуть перевищувати 60. 
Набрані на екзамені бали не можуть перевищувати 40. 
Набрані на екзамені бали повинні бути не менше 24 або дорівнювати 0(в останньому випадку оцінка не може бути задовільною). 
Сумарна оцінка в балах є сумою балів, набраних у семестрі та на екзамені.
Сумарна оцінка в балах та за державною шкалою мають бути узгоджені між собою.

Прізвище, ім’я та по-батькові (якщо наявне у варіанті) можуть мати до 26, 28, 27 відповідно символів включно кожне. 
Крім літер алфавіту можуть містити апостроф, дефіс та пробіл.

Предмет може мати від 4 до 66 символів включно. 
Крім літер алфавіту може містити подвійні лапки, пробіл та дефіс.

Код групи є непорожнім та складається не більш ніж з 7 літер. 
Крім літер алфавіту може містити десяткові цифри та дефіс.


Номер залікової книжки однозначно ідентифікує студента та складається з 8 десяткових цифр. 

\bigskip\textit{Для деяких варіантів може знадобитись}

Середній бал/оцінка за предметом визначається як середнє арифметичнесумарних балів/державних оцінок за цим предметом. 
Середній бал/оцінка студента визначається як середнє арифметичне сумарних балів/державних оцінокцього студента. 
За обчислення середньої оцінки недопущені та ті, що не з’явилися, рахуються як 2.
У решті випадків недопуски підсумовуються як $-1$, а «не з’явився» як 0.
Якщо студент не гірше ніж задовільно склав усі предмети, 
то його рейтинг за балами/оцінкою визначається як середній бал/оцінка; інакше дорівнює 0.

Стабільність студента визначається як модуль різниці між його кращим та найгіршим сумарним балом.
Стабільність предмета визначається як модуль різниці між найкращим та найгіршимсумарним балом за цим предметом.




\par
\newpage
{\center Варіант  \No \textbf { 25 }\par}
\bigskip
%type = session
%variant = 56
Обробляються результати складання студентами екзаменів зимової сесії.

\underline{Записи основного файлу містять поля}: прізвище, предмет, код групи, по-батькові, ім’я, оцінка за державною шкалою, номер залікової книжки, сумарна оцінка в балах за 100-бальною системою, набрані на екзамені бали.

\underline{У допоміжному файлі наявні ключі}: кількість оцінок «відмінно», кількість оцінок «задовільно», загальна кількість записів.

\underline{Знайти} всіх студентів, що щось не склали. Вивести по кожному з них інформацію:\par
{\noindent}--- на першому рядку:\par
середній бал (округлений до цілого), кількість хвостів, номер залікової книжки, прізвище, ім’я, по-батькові;

{\noindent}--- на наступних рядках, починаючи з табуляції, вивести для студента всі його результати (по одному на рядок):\par
предмет, державна оцінка прописом

{\noindent}у такому сортуванні: державна оцінка прописом (впорядковуємо як рядки), предмет.

%Студентів сортувати: кількість хвостів (за спаданням), середній бал (після округлення), номер залікової книжки.

\bigskip\textit{Обмеження}

Студент не може складати предмет більше одного разу.

Бали та оцінка є значеннями цілих числових типів. Решта полів є непорожніми рядками.
Рядки розпочинатись та закінчуватись білими символами не можуть.

Оцінка за державною шкалою може приймати значення: 2-5~--- як зазвичай, 0~--- не з’явився,$-1$~--- не допущений. 
Оцінки 3-5 вважаються задовільними, решта~--- незадовільними. 
Бали є цілими та невід’ємними.
Набрані впродовж семестру бали не можуть перевищувати 60. 
Набрані на екзамені бали не можуть перевищувати 40. 
Набрані на екзамені бали повинні бути не менше 24 або дорівнювати 0(в останньому випадку оцінка не може бути задовільною). 
Сумарна оцінка в балах є сумою балів, набраних у семестрі та на екзамені.
Сумарна оцінка в балах та за державною шкалою мають бути узгоджені між собою.

Прізвище, ім’я та по-батькові (якщо наявне у варіанті) можуть мати до 23, 30, 26 відповідно символів включно кожне. 
Крім літер алфавіту можуть містити апостроф, дефіс та пробіл.

Предмет може мати від 3 до 50 символів включно. 
Крім літер алфавіту може містити подвійні лапки, пробіл та дефіс.

Код групи є непорожнім та складається не більш ніж з 3 літер. 
Крім літер алфавіту може містити десяткові цифри та дефіс.


Номер залікової книжки однозначно ідентифікує студента та складається з 6 десяткових цифр. 

\bigskip\textit{Для деяких варіантів може знадобитись}

Середній бал/оцінка за предметом визначається як середнє арифметичнесумарних балів/державних оцінок за цим предметом. 
Середній бал/оцінка студента визначається як середнє арифметичне сумарних балів/державних оцінокцього студента. 
За обчислення середньої оцінки недопущені та ті, що не з’явилися, рахуються як 2.
У решті випадків недопуски підсумовуються як $-1$, а «не з’явився» як 0.
Якщо студент не гірше ніж задовільно склав усі предмети, 
то його рейтинг за балами/оцінкою визначається як середній бал/оцінка; інакше дорівнює 0.

Стабільність студента визначається як модуль різниці між його кращим та найгіршим сумарним балом.
Стабільність предмета визначається як модуль різниці між найкращим та найгіршимсумарним балом за цим предметом.




\par
\newpage
{\center Варіант  \No \textbf { 26 }\par}
\bigskip
%type = session
%variant = 77
Обробляються результати складання студентами екзаменів зимової сесії.

\underline{Записи основного файлу містять поля}: прізвище, по-батькові, предмет, набрані впродовж семестру бали, номер залікової книжки, ім’я, набрані на екзамені бали, оцінка за державною шкалою, код групи, сумарна оцінка в балах за 100-бальною системою.

\underline{У допоміжному файлі наявні ключі}: сума державних оцінок, загальна кількість записів.

\underline{Знайти} предмети, за якими не було жодної оцінки «добре». Вивести по кожному з них інформацію:\par
{\noindent}--- на першому рядку:\par
предмет, кількість оцінок «відмінно», середній бал за предметом (округлений до цілого);

{\noindent}--- на наступних рядках, починаючи з табуляції, вивести для предмета всіх студентів, що його складали (по одному на рядок):\par
група, кількість оцінок «відмінно», середній бал по групі (округлений до цілого)

{\noindent}у такому сортуванні: група.

%Предмети сортувати: середній бал (округлений), предмет.

\bigskip\textit{Обмеження}

Студент не може складати предмет більше одного разу.

Бали та оцінка є значеннями цілих числових типів. Решта полів є непорожніми рядками.
Рядки розпочинатись та закінчуватись білими символами не можуть.

Оцінка за державною шкалою може приймати значення: 2-5~--- як зазвичай, 0~--- не з’явився,$-1$~--- не допущений. 
Оцінки 3-5 вважаються задовільними, решта~--- незадовільними. 
Бали є цілими та невід’ємними.
Набрані впродовж семестру бали не можуть перевищувати 60. 
Набрані на екзамені бали не можуть перевищувати 40. 
Набрані на екзамені бали повинні бути не менше 24 або дорівнювати 0(в останньому випадку оцінка не може бути задовільною). 
Сумарна оцінка в балах є сумою балів, набраних у семестрі та на екзамені.
Сумарна оцінка в балах та за державною шкалою мають бути узгоджені між собою.

Прізвище, ім’я та по-батькові (якщо наявне у варіанті) можуть мати до 29, 20, 25 відповідно символів включно кожне. 
Крім літер алфавіту можуть містити апостроф, дефіс та пробіл.

Предмет може мати від 5 до 51 символів включно. 
Крім літер алфавіту може містити подвійні лапки, пробіл та дефіс.

Код групи є непорожнім та складається не більш ніж з 3 літер. 
Крім літер алфавіту може містити десяткові цифри та дефіс.


Номер залікової книжки однозначно ідентифікує студента та складається з 8 десяткових цифр. 

\bigskip\textit{Для деяких варіантів може знадобитись}

Середній бал/оцінка за предметом визначається як середнє арифметичнесумарних балів/державних оцінок за цим предметом. 
Середній бал/оцінка студента визначається як середнє арифметичне сумарних балів/державних оцінокцього студента. 
За обчислення середньої оцінки недопущені та ті, що не з’явилися, рахуються як 2.
У решті випадків недопуски підсумовуються як $-1$, а «не з’явився» як 0.
Якщо студент не гірше ніж задовільно склав усі предмети, 
то його рейтинг за балами/оцінкою визначається як середній бал/оцінка; інакше дорівнює 0.

Стабільність студента визначається як модуль різниці між його кращим та найгіршим сумарним балом.
Стабільність предмета визначається як модуль різниці між найкращим та найгіршимсумарним балом за цим предметом.




\par
\newpage
{\center Варіант  \No \textbf { 27 }\par}
\bigskip
%type = session
%variant = 61
Обробляються результати складання студентами екзаменів зимової сесії.

\underline{Записи основного файлу містять поля}: по-батькові, сумарна оцінка в балах за 100-бальною системою, набрані впродовж семестру бали, предмет, прізвище, ім’я, код групи, номер залікової книжки, набрані на екзамені бали, оцінка за державною шкалою.

\underline{У допоміжному файлі наявні ключі}: кількість «незадовільно», кількість сумарних оцінок від 60 до 70 включно.

\underline{Знайти} всіх студентів, що по кожному предмету отримали оцінку 3 за державною шкалою. Вивести по кожному з них інформацію:\par
{\noindent}--- на першому рядку:\par
середня оцінка в балах (округлена до цілого), прізвище, ім’я, номер залікової книжки, кількість екзаменів;

{\noindent}--- на наступних рядках, починаючи з табуляції, вивести для студента всі його результати (по одному на рядок):\par
предмет, сумарна оцінка, бали на екзамені, бали в семестрі

{\noindent}у такому сортуванні: бали на екзамені, предмет.

%Студентів сортувати: середня оцінка в балах (після округлення), прізвище, ім’я, номер залікової книжки.

\bigskip\textit{Обмеження}

Студент не може складати предмет більше одного разу.

Бали та оцінка є значеннями цілих числових типів. Решта полів є непорожніми рядками.
Рядки розпочинатись та закінчуватись білими символами не можуть.

Оцінка за державною шкалою може приймати значення: 2-5~--- як зазвичай, 0~--- не з’явився,$-1$~--- не допущений. 
Оцінки 3-5 вважаються задовільними, решта~--- незадовільними. 
Бали є цілими та невід’ємними.
Набрані впродовж семестру бали не можуть перевищувати 60. 
Набрані на екзамені бали не можуть перевищувати 40. 
Набрані на екзамені бали повинні бути не менше 24 або дорівнювати 0(в останньому випадку оцінка не може бути задовільною). 
Сумарна оцінка в балах є сумою балів, набраних у семестрі та на екзамені.
Сумарна оцінка в балах та за державною шкалою мають бути узгоджені між собою.

Прізвище, ім’я та по-батькові (якщо наявне у варіанті) можуть мати до 30, 23, 21 відповідно символів включно кожне. 
Крім літер алфавіту можуть містити апостроф, дефіс та пробіл.

Предмет може мати від 4 до 56 символів включно. 
Крім літер алфавіту може містити подвійні лапки, пробіл та дефіс.

Код групи є непорожнім та складається не більш ніж з 5 літер. 
Крім літер алфавіту може містити десяткові цифри та дефіс.


Номер залікової книжки однозначно ідентифікує студента та складається з 7 десяткових цифр. 

\bigskip\textit{Для деяких варіантів може знадобитись}

Середній бал/оцінка за предметом визначається як середнє арифметичнесумарних балів/державних оцінок за цим предметом. 
Середній бал/оцінка студента визначається як середнє арифметичне сумарних балів/державних оцінокцього студента. 
За обчислення середньої оцінки недопущені та ті, що не з’явилися, рахуються як 2.
У решті випадків недопуски підсумовуються як $-1$, а «не з’явився» як 0.
Якщо студент не гірше ніж задовільно склав усі предмети, 
то його рейтинг за балами/оцінкою визначається як середній бал/оцінка; інакше дорівнює 0.

Стабільність студента визначається як модуль різниці між його кращим та найгіршим сумарним балом.
Стабільність предмета визначається як модуль різниці між найкращим та найгіршимсумарним балом за цим предметом.




\par
\newpage
{\center Варіант  \No \textbf { 28 }\par}
\bigskip
%type = session
%variant = 76
Обробляються результати складання студентами екзаменів зимової сесії.

\underline{Записи основного файлу містять поля}: набрані на екзамені бали, оцінка за державною шкалою, ім’я, номер залікової книжки, по-батькові, сумарна оцінка в балах за 100-бальною системою, код групи, прізвище, предмет.

\underline{У допоміжному файлі наявні ключі}: кількість оцінок «добре», загальна кількість записів.

\underline{Знайти} предмети, стабільність яких вища середньої стабільності по всіх предметах. Вивести по кожному з них інформацію:\par
{\noindent}--- на першому рядку:\par
стабільність, предмет, кількість студентів, що не склали;

{\noindent}--- на наступних рядках, починаючи з табуляції, вивести для предмета всіх студентів, що його складали (по одному на рядок):\par
бали за екзамен, оцінка в балах, ім’я, прізвище, номер залікової книжки, група

{\noindent}у такому сортуванні: група, номер залікової книжки.

%Предмети сортувати: стабільність (за спаданням), кількість студентів, що не склали, предмет.

\bigskip\textit{Обмеження}

Студент не може складати предмет більше одного разу.

Бали та оцінка є значеннями цілих числових типів. Решта полів є непорожніми рядками.
Рядки розпочинатись та закінчуватись білими символами не можуть.

Оцінка за державною шкалою може приймати значення: 2-5~--- як зазвичай, 0~--- не з’явився,$-1$~--- не допущений. 
Оцінки 3-5 вважаються задовільними, решта~--- незадовільними. 
Бали є цілими та невід’ємними.
Набрані впродовж семестру бали не можуть перевищувати 60. 
Набрані на екзамені бали не можуть перевищувати 40. 
Набрані на екзамені бали повинні бути не менше 24 або дорівнювати 0(в останньому випадку оцінка не може бути задовільною). 
Сумарна оцінка в балах є сумою балів, набраних у семестрі та на екзамені.
Сумарна оцінка в балах та за державною шкалою мають бути узгоджені між собою.

Прізвище, ім’я та по-батькові (якщо наявне у варіанті) можуть мати до 20, 24, 24 відповідно символів включно кожне. 
Крім літер алфавіту можуть містити апостроф, дефіс та пробіл.

Предмет може мати від 3 до 52 символів включно. 
Крім літер алфавіту може містити подвійні лапки, пробіл та дефіс.

Код групи є непорожнім та складається не більш ніж з 6 літер. 
Крім літер алфавіту може містити десяткові цифри та дефіс.


Номер залікової книжки однозначно ідентифікує студента та складається з 7 десяткових цифр. 

\bigskip\textit{Для деяких варіантів може знадобитись}

Середній бал/оцінка за предметом визначається як середнє арифметичнесумарних балів/державних оцінок за цим предметом. 
Середній бал/оцінка студента визначається як середнє арифметичне сумарних балів/державних оцінокцього студента. 
За обчислення середньої оцінки недопущені та ті, що не з’явилися, рахуються як 2.
У решті випадків недопуски підсумовуються як $-1$, а «не з’явився» як 0.
Якщо студент не гірше ніж задовільно склав усі предмети, 
то його рейтинг за балами/оцінкою визначається як середній бал/оцінка; інакше дорівнює 0.

Стабільність студента визначається як модуль різниці між його кращим та найгіршим сумарним балом.
Стабільність предмета визначається як модуль різниці між найкращим та найгіршимсумарним балом за цим предметом.




\par
\newpage
{\center Варіант  \No \textbf { 29 }\par}
\bigskip
%type = session
%variant = 64
Обробляються результати складання студентами екзаменів зимової сесії.

\underline{Записи основного файлу містять поля}: оцінка за державною шкалою, набрані на екзамені бали, предмет, код групи, набрані впродовж семестру бали, ім’я, номер залікової книжки, сумарна оцінка в балах за 100-бальною системою, прізвище.

\underline{У допоміжному файлі наявні ключі}: найбільший сумарний бал, найбільший бал в семестрі, кількість студентів.

\underline{Знайти} всіх студентів, чий рейтинг за балами вище середнього рейтингу. Вивести по кожному з них інформацію:\par
{\noindent}--- на першому рядку:\par
ім’я, прізвище, рейтинг за балами (округлений до цілого), номер залікової книжки, кількість «відмінно»;

{\noindent}--- на наступних рядках, починаючи з табуляції, вивести для студента всі його результати (по одному на рядок):\par
державна оцінка, предмет, сумарна оцінка

{\noindent}у такому сортуванні: державна оцінка, предмет.

%Студентів сортувати: ім’я, прізвище, рейтинг за балами (після округлення), номер залікової книжки.

\bigskip\textit{Обмеження}

Студент не може складати предмет більше одного разу.

Бали та оцінка є значеннями цілих числових типів. Решта полів є непорожніми рядками.
Рядки розпочинатись та закінчуватись білими символами не можуть.

Оцінка за державною шкалою може приймати значення: 2-5~--- як зазвичай, 0~--- не з’явився,$-1$~--- не допущений. 
Оцінки 3-5 вважаються задовільними, решта~--- незадовільними. 
Бали є цілими та невід’ємними.
Набрані впродовж семестру бали не можуть перевищувати 60. 
Набрані на екзамені бали не можуть перевищувати 40. 
Набрані на екзамені бали повинні бути не менше 24 або дорівнювати 0(в останньому випадку оцінка не може бути задовільною). 
Сумарна оцінка в балах є сумою балів, набраних у семестрі та на екзамені.
Сумарна оцінка в балах та за державною шкалою мають бути узгоджені між собою.

Прізвище, ім’я та по-батькові (якщо наявне у варіанті) можуть мати до 21, 29, 20 відповідно символів включно кожне. 
Крім літер алфавіту можуть містити апостроф, дефіс та пробіл.

Предмет може мати від 4 до 67 символів включно. 
Крім літер алфавіту може містити подвійні лапки, пробіл та дефіс.

Код групи є непорожнім та складається не більш ніж з 7 літер. 
Крім літер алфавіту може містити десяткові цифри та дефіс.


Номер залікової книжки однозначно ідентифікує студента та складається з 6 десяткових цифр. 

\bigskip\textit{Для деяких варіантів може знадобитись}

Середній бал/оцінка за предметом визначається як середнє арифметичнесумарних балів/державних оцінок за цим предметом. 
Середній бал/оцінка студента визначається як середнє арифметичне сумарних балів/державних оцінокцього студента. 
За обчислення середньої оцінки недопущені та ті, що не з’явилися, рахуються як 2.
У решті випадків недопуски підсумовуються як $-1$, а «не з’явився» як 0.
Якщо студент не гірше ніж задовільно склав усі предмети, 
то його рейтинг за балами/оцінкою визначається як середній бал/оцінка; інакше дорівнює 0.

Стабільність студента визначається як модуль різниці між його кращим та найгіршим сумарним балом.
Стабільність предмета визначається як модуль різниці між найкращим та найгіршимсумарним балом за цим предметом.




\par
\newpage
{\center Варіант  \No \textbf { 30 }\par}
\bigskip
%type = session
%variant = 74
Обробляються результати складання студентами екзаменів зимової сесії.

\underline{Записи основного файлу містять поля}: оцінка за державною шкалою, набрані на екзамені бали, номер залікової книжки, прізвище, набрані впродовж семестру бали, сумарна оцінка в балах за 100-бальною системою, ім’я, предмет, код групи.

\underline{У допоміжному файлі наявні ключі}: найменша довжина назви предмета, загальна кількість записів.

\underline{Знайти} предмети, середній бал за які, вищий середнього за всю сесію. Вивести по кожному з них інформацію:\par
{\noindent}--- на першому рядку:\par
предмет, середній бал за предметом (округлений до 1 знаку), кількість студентів, що склали на «відмінно»;

{\noindent}--- на наступних рядках, починаючи з табуляції, вивести для предмета студентів, що його не склали (по одному на рядок):\par
прізвище, ім’я, по-батькові, номер залікової книжки, група, оцінка за державною шкалою прописом

{\noindent}у такому сортуванні: номер залікової книжки.

%Предмети сортувати: середній бал (округлений), кількість студентів, що склали на «відмінно», предмет.

\bigskip\textit{Обмеження}

Студент не може складати предмет більше одного разу.

Бали та оцінка є значеннями цілих числових типів. Решта полів є непорожніми рядками.
Рядки розпочинатись та закінчуватись білими символами не можуть.

Оцінка за державною шкалою може приймати значення: 2-5~--- як зазвичай, 0~--- не з’явився,$-1$~--- не допущений. 
Оцінки 3-5 вважаються задовільними, решта~--- незадовільними. 
Бали є цілими та невід’ємними.
Набрані впродовж семестру бали не можуть перевищувати 60. 
Набрані на екзамені бали не можуть перевищувати 40. 
Набрані на екзамені бали повинні бути не менше 24 або дорівнювати 0(в останньому випадку оцінка не може бути задовільною). 
Сумарна оцінка в балах є сумою балів, набраних у семестрі та на екзамені.
Сумарна оцінка в балах та за державною шкалою мають бути узгоджені між собою.

Прізвище, ім’я та по-батькові (якщо наявне у варіанті) можуть мати до 25, 25, 22 відповідно символів включно кожне. 
Крім літер алфавіту можуть містити апостроф, дефіс та пробіл.

Предмет може мати від 5 до 54 символів включно. 
Крім літер алфавіту може містити подвійні лапки, пробіл та дефіс.

Код групи є непорожнім та складається не більш ніж з 4 літер. 
Крім літер алфавіту може містити десяткові цифри та дефіс.


Номер залікової книжки однозначно ідентифікує студента та складається з 8 десяткових цифр. 

\bigskip\textit{Для деяких варіантів може знадобитись}

Середній бал/оцінка за предметом визначається як середнє арифметичнесумарних балів/державних оцінок за цим предметом. 
Середній бал/оцінка студента визначається як середнє арифметичне сумарних балів/державних оцінокцього студента. 
За обчислення середньої оцінки недопущені та ті, що не з’явилися, рахуються як 2.
У решті випадків недопуски підсумовуються як $-1$, а «не з’явився» як 0.
Якщо студент не гірше ніж задовільно склав усі предмети, 
то його рейтинг за балами/оцінкою визначається як середній бал/оцінка; інакше дорівнює 0.

Стабільність студента визначається як модуль різниці між його кращим та найгіршим сумарним балом.
Стабільність предмета визначається як модуль різниці між найкращим та найгіршимсумарним балом за цим предметом.




\par
\newpage
\end{document}
